\chapter{质点动力学}

\section{牛顿运动定律}

\begin{problem}{2.1.1}
质量为 $M$ 的人通过质量可以不计、不可伸长的绳子和摩擦可忽略的定滑轮，

将质量为 m 的重物提起, 如图 2.1 所示. 分两种情况: (1) 物

体匀速上升；(2)物体以加速度 $a$ 上升.求人对地面的压力.
\begin{center}
\begin{tikzpicture}[>=Stealth, every node/.style={font=\small}]
  % Ceiling
  \draw[thick] (-1.5, 4.2) -- (1.5, 4.2);
  \fill[pattern=north east lines] (-1.5, 4.2) rectangle (1.5, 4.45);

  % Pulley support & pulley
  \draw[thick] (0, 4.2) -- (0, 2.9);
  \draw[thick] (0, 2.9) circle (0.95);
  \fill (0, 2.9) circle (1.2pt);

  % Left rope and person M
  \draw[thick] (-0.95, 2.9) -- (-0.95, 0.55);
  % Stick figure
  \draw[thick] (-1.5, 1.1) circle (0.17);
  \draw[thick] (-1.45, 0.93) -- (-1.45, 0.42);
  \draw[thick] (-1.45, 0.42) -- (-1.7, 0);
  \draw[thick] (-1.45, 0.42) -- (-1.2, 0);
  \draw[thick] (-1.45, 0.85) .. controls (-1.25, 0.72) .. (-0.95, 0.98);
  \draw[thick] (-1.45, 0.72) .. controls (-1.25, 0.60) .. (-0.95, 0.90);
  \node[left] at (-1.65, 0.65) {$M$};

  % Right rope and hanging block m
  \draw[thick] (0.95, 2.9) -- (0.95, 1.3);
  \draw[thick] (0.55, 0.65) rectangle (1.35, 1.3);
  \node[right] at (1.4, 0.95) {$m$};

  % Ground
  \draw[thick] (-2.1, 0) -- (2.1, 0);
  \fill[pattern=north east lines] (-2.1, -0.25) rectangle (2.1, 0);
\end{tikzpicture}
\end{center}
\end{problem}
\clearpage


\begin{problem}{2.1.2}
图 2.2 中人的质量 $m_{1}=60kg$ ，人所站着的底板质量 $m_{2}=20kg$ ，绳子和滑轮质量均可忽略不计。人要用多大的力拉住绳子，才能保持自己和底板静止不动。
\begin{center}
\begin{tikzpicture}[>=Stealth, scale=0.9, every node/.style={font=\small}]
  % Ceiling
  \draw[thick] (-1.5, 5.5) -- (2.6, 5.5);
  \fill[pattern=north east lines] (-1.5, 5.5) rectangle (2.6, 5.75);

  % Pulley 1 (top)
  \draw[thick] (0, 5.5) -- (0, 4.35);
  \draw[thick] (0, 4.35) circle (0.75);
  \draw[thick] (-0.75, 4.35) -- (-0.75, 0.3);
  \draw[thick] (0.75, 4.35) -- (0.75, 2.95);

  % Pulley 2 (middle)
  \draw[thick] (0.75, 2.95) circle (0.75);
  \draw[thick] (0, 2.95) -- (0, 0.3);
  \draw[thick] (1.5, 2.95) -- (1.5, 1.55);

  % Pulley 3 (bottom)
  \draw[thick] (1.5, 1.55) circle (0.75);
  \draw[thick] (0.75, 1.55) -- (0.75, 0.3);
  \draw[thick] (2.25, 1.55) -- (2.25, 0.3);

  % Person standing on platform holding left rope of Pulley 3
  \draw[thick] (0.3, 1.15) circle (0.14);
  \draw[thick] (0.35, 1.01) -- (0.35, 0.58);
  \draw[thick] (0.35, 0.58) -- (0.18, 0.3);
  \draw[thick] (0.35, 0.58) -- (0.50, 0.3);
  \draw[thick] (0.35, 0.92) -- (0.55, 0.78) -- (0.75, 0.98);
  \draw[thick] (0.35, 0.80) -- (0.55, 0.68) -- (0.75, 0.90);
  \node[right] at (0.8, 0.6) {$m_1$};

  % Platform
  \draw[thick] (-1.6, 0.05) rectangle (2.9, 0.3);
  \node[below] at (0.35, 0.02) {$m_2$};
\end{tikzpicture}
\end{center}
\end{problem}
\clearpage


\begin{problem}{2.1.3}
跨过一个无摩擦定滑轮的一根绳子,一端挂有9kg质量,另一端挂有7kg质量,如图2.3所示.求加速度和绳中的张力.

解得
\begin{center}
\begin{tikzpicture}[>=Stealth, every node/.style={font=\small}]
  % Pulley
  \draw[thick] (0, 3.0) circle (1.1);
  \draw[thick] (0, 3.0) circle (0.08);

  % Left rope and m1
  \draw[thick] (-1.1, 3.0) -- (-1.1, 1.35);
  \draw[thick] (-1.5, 0.85) rectangle (-0.7, 1.35);
  \node[right] at (-0.65, 1.05) {$m_1$};

  % Right rope and m2
  \draw[thick] (1.1, 3.0) -- (1.1, 0.6);
  \draw[thick] (0.7, 0.1) rectangle (1.5, 0.6);
  \node[right] at (1.55, 0.3) {$m_2$};
\end{tikzpicture}
\end{center}
\end{problem}
\clearpage


\begin{problem}{2.1.4}
图 2.4 中滑轮及绳子质量均可忽略不计, 所有的接触均是光滑的. 分两种情况: (1) 要使 $m_{1}, m_{2}, M$ 之间无相对运动; (2) 要使 $M$ 保持静止, 必须施以 $M$ 的水平力 $F$ 多大? $M$ 对水平面的压力 $N$ 多大?
\begin{center}
\begin{tikzpicture}[>=Stealth, every node/.style={font=\small}]
  % Ground
  \draw[thick] (-1.1, 0) -- (4.9, 0);
  \fill[pattern=north east lines] (-1.1, -0.25) rectangle (4.9, 0);

  % Block M
  \draw[thick] (0, 0) rectangle (3.4, 2.2);
  \node at (1.7, 0.85) {$M$};

  % Force F
  \draw[->, thick] (-1.0, 1.1) -- node[above] {$F$} (0, 1.1);

  % Block m1
  \draw[thick] (0.6, 2.2) rectangle (1.4, 3.5);
  \node at (1.0, 2.85) {$m_1$};

  % Pulley at top-right of M
  \coordinate (P) at (3.85, 2.65);
  \draw[thick] (3.4, 2.2) -- (P);
  \draw[thick, fill=white] (P) circle (0.45);
  \draw[thick] (P) circle (0.06);

  % Rope
  \draw[thick] (1.4, 3.10) -- (3.85, 3.10);
  \draw[thick] (4.30, 2.65) -- (4.30, 1.65);

  % Block m2
  \draw[thick] (3.4, 0.95) rectangle (4.55, 1.65);
  \node at (3.98, 1.3) {$m_2$};
\end{tikzpicture}
\end{center}
\end{problem}
\clearpage


\begin{problem}{2.1.5}
图 2.5 中小车质量 M、重物质量 m 以及固定斜面的倾角 $\alpha$ 都是已知的，绳子不可伸长，绳子和滑轮的质量以及整个装置中的摩擦力均可略去不计，求小车和重物的加速度.
\begin{center}
\begin{tikzpicture}[>=Stealth, scale=0.92, every node/.style={font=\small}]
  % Inclined plane angle alpha = 25 deg
  \def\ang{25}
  \coordinate (R) at (5.6, 0);
  \coordinate (T) at ($(R) + (180-\ang:4.2)$); % top of slope (~1.79, 1.775)
  \coordinate (Vbase) at (1.79, 0);

  % Ground
  \draw[thick] (0.8, 0) -- (5.8, 0);
  \fill[pattern=north east lines] (0.8, -0.25) rectangle (5.8, 0);

  % Incline & support stand
  \draw[thick] (R) -- (T) -- (Vbase);
  % Perpendicular stop plate at top of incline
  \coordinate (StopTop) at ($(T) + (90-\ang:0.55)$);
  \coordinate (StopBot) at ($(T) + (-90-\ang:0.15)$);
  \draw[thick] (StopBot) -- (StopTop);

  % Pulley 2 center and mounting bracket
  \coordinate (P2) at (1.40, 2.65);
  \draw[thick] (T) -- (P2) -- ($(T) + (90-\ang:0.35)$);
  \draw[thick, fill=white] (P2) circle (0.55);
  \node at ($(P2) + (90:0.78)$) {$2$};

  % Cart M on incline
  % Along slope from R: distance 1.45 to center of cart
  \coordinate (CartBase) at ($(R) + (180-\ang:1.25)$);
  \begin{scope}[shift={(CartBase)}, rotate=-\ang]
    % Two wheels
    \draw[thick, fill=white] (-0.38, 0.14) circle (0.14);
    \draw[thick, fill=white] (0.38, 0.14) circle (0.14);
    % Cart body
    \draw[thick, fill=white] (-0.65, 0.28) rectangle (0.65, 1.0);
    \node at (0, 0.64) {$M$};
  \end{scope}

  % Pulley 1 (attached to cart M)
  \coordinate (CartFront) at ($(CartBase) + (180-\ang:0.65) + (90-\ang:0.64)$);
  \coordinate (P1) at ($(CartFront) + (180-\ang:0.95)$);
  \draw[thick] (CartFront) -- (P1);
  \draw[thick, fill=white] (P1) circle (0.52);
  \node at ($(P1) + (90:0.75)$) {$1$};

  % Ropes for Pulley 1 and Pulley 2
  % Bottom tangent of P1 to stop plate
  \coordinate (P1bot) at ($(P1) + (-90-\ang:0.52)$);
  \coordinate (StopAnchor) at ($(T) + (90-\ang:0.25)$);
  \draw[thick] (P1bot) -- (StopAnchor);
  % Top tangent of P1 to top-right tangent of P2
  \draw[thick] ($(P1) + (90-\ang:0.52)$) -- ($(P2) + (90-\ang:0.55)$);

  % Pulley 3 (movable pulley)
  \coordinate (P3) at (0.85, 1.25);
  \draw[thick] (0.85, 2.65) -- (P3); % rope from left of P2 to center of P3
  \draw[thick, fill=white] (P3) circle (0.52);
  \node at ($(P3) + (135:0.75)$) {$3$};

  % Rope over Pulley 3: right side anchored to ground, left side to mass m
  \draw[thick] (1.37, 1.25) -- (1.37, 0);
  \draw[thick] (1.37, 0.25) -- (1.79, 0.25); % base brace to stand
  \draw[thick] (0.33, 1.25) -- (0.33, -0.35);

  % Hanging mass m
  \draw[thick] (0.0, -0.85) rectangle (0.66, -0.35);
  \node at (0.33, -0.60) {$m$};

  % Angle alpha
  \draw (4.7, 0) arc (180:180-\ang:0.9);
  \node at (4.55, 0.21) {$\alpha$};
\end{tikzpicture}
\end{center}
\end{problem}
\clearpage


\begin{problem}{2.1.6}
三块质量均为 $m$ 的相同物块叠放在水平面上, 如图 2.6 所示. 各接触面间的摩擦因数均为 $\mu$ , 有一从零不断加大的水平力 $F$ 作用在最下面的物块上. 问 $F$ 达多大时, 最下面的物块才能有对上面两物块的相对运动.
\begin{center}
\begin{tikzpicture}[>=Stealth, every node/.style={font=\small}]
  % Ground
  \draw[thick] (-1.8, 0) -- (2.1, 0);
  \fill[pattern=north east lines] (-1.8, -0.25) rectangle (2.1, 0);

  % Three stacked blocks of mass m
  \draw[thick] (-0.95, 0) rectangle (0.95, 0.75);
  \node at (0, 0.375) {$m$};
  \draw[thick] (-0.95, 0.75) rectangle (0.95, 1.5);
  \node at (0, 1.125) {$m$};
  \draw[thick] (-0.95, 1.5) rectangle (0.95, 2.25);
  \node at (0, 1.875) {$m$};

  % Force F on bottom block
  \draw[->, thick] (0.95, 0.375) -- (2.05, 0.375) node[right] {$F$};
\end{tikzpicture}
\end{center}
\end{problem}
\clearpage


\begin{problem}{2.1.7}
质量为 $M$ 的木板静置于水平桌面上(如图2.7所示)，一端 $A$ 与桌边对齐，木板上放着一质量为 $m$ 的小物体，它离板的 $A$ 端距离 $l$ ，桌面长 $L$ .现在板的另一端 $B$ 有一恒定的水平力 $F$ 作用，要将木板从小物体下抽出，又不至于使小物体从桌边落到地上.若各接触面的摩擦因数均为 $\mu ,F$ 至少多大？
\begin{center}
\begin{tikzpicture}[>=Stealth, scale=0.92, every node/.style={font=\small}]
  % Ground
  \draw[thick] (-0.4, 0) -- (4.8, 0);
  \fill[pattern=north east lines] (-0.4, -0.25) rectangle (4.8, 0);

  % Table legs
  \draw[thick] (0.65, 1.1) -- (0.4, 0) -- (0.6, 0) -- (0.85, 1.1);
  \draw[thick] (2.95, 1.1) -- (3.2, 0) -- (3.4, 0) -- (3.15, 1.1);

  % Tabletop (length L = 3.6)
  \draw[thick] (0, 1.1) rectangle (3.6, 1.45);

  % Plank M (from A at x=0 to B at x=4.8)
  \draw[thick] (0, 1.45) rectangle (4.8, 1.72);
  \node[above] at (0, 1.75) {$A$};
  \node[above] at (4.8, 1.75) {$B$};
  \node[above] at (3.35, 1.75) {$M$};

  % Block m on plank (centered at x=1.7)
  \draw[thick] (1.3, 1.72) rectangle (2.1, 2.25);
  \node at (1.7, 1.98) {$m$};

  % Force F
  \draw[->, thick] (4.8, 1.585) -- (5.6, 1.585) node[right] {$F$};

  % Dimension l (from A to center of m)
  \draw (0, 2.35) -- (0, 2.85);
  \draw (1.7, 2.35) -- (1.7, 2.85);
  \draw[<->] (0, 2.65) -- node[above] {$l$} (1.7, 2.65);

  % Dimension L (tabletop length)
  \draw (0, 0.45) -- (0, 0.95);
  \draw (3.6, 0.45) -- (3.6, 0.95);
  \draw[<->] (0, 0.65) -- node[above] {$L$} (3.6, 0.65);
\end{tikzpicture}
\end{center}
\end{problem}
\clearpage


\begin{problem}{2.1.8}
一块砖以初速度 $1.5\mathrm{m / s}$ 在一与水平面成 $30^{\circ}$ 角的斜面上向上运动，摩擦因数 $\mu = \frac{\sqrt{3}}{12}$ . 问经 $0.5\mathrm{s}$ 后，此砖离它的初始位置有多远？
\end{problem}
\clearpage


\begin{problem}{2.1.9}
两个无质量的小环在一个处于铅垂面内的光滑圆环上滑动,一穿过两环的光滑弦线上带有三个重物,两个在两端,一个在两环之间.如果当小环 $A$ 、 $B$ 在离开圆环最高点 $H$ 对环心 $O$ 的张角 $30^{\circ}$ 时处于平衡,如图2.8所示.求系在弦线上的三重物 $D$ 、 $E$ 、 $G$ 的质量之间的关系.
\begin{center}
\begin{tikzpicture}[>=Stealth, every node/.style={font=\small}]
  \def\R{2.3}
  \coordinate (O) at (0, 0);
  \coordinate (H) at (0, \R);
  \coordinate (A) at (120:\R);
  \coordinate (B) at (60:\R);
  \coordinate (G) at (0, 1.50);
  \coordinate (D) at ($(A) + (0, -1.4)$);
  \coordinate (E) at ($(B) + (0, -1.65)$);

  % Vertical circular hoop
  \draw[thick] (O) circle (\R);

  % Dashed radii OH, OA, OB
  \draw[dashed] (O) -- (H);
  \draw[dashed] (O) -- (A);
  \draw[dashed] (O) -- (B);

  % Angles 30 deg
  \draw (0, 0.85) arc (90:120:0.85);
  \node at (107:1.15) {$30^\circ$};
  \draw (0, 0.85) arc (90:60:0.85);
  \node at (73:1.15) {$30^\circ$};

  % Small rings A and B
  \draw[thick] (A) circle (0.13);
  \draw[thick] (B) circle (0.13);

  % String A-G-B and vertical ends to D and E
  \draw[thick] (A) -- (G) -- (B);
  \draw[thick] (A) -- (D);
  \draw[thick] (B) -- (E);

  % Weights G, D, E
  \fill (G) circle (2pt);
  \draw[thick, fill=white] (D) circle (2pt);
  \draw[thick, fill=white] (E) circle (2pt);

  % Labels
  \node[above] at (H) {$H$};
  \node[above left] at ($(A)+(-0.08,0.05)$) {$A$};
  \node[above right] at ($(B)+(0.08,0.05)$) {$B$};
  \node at (0.28, 1.85) {$G$};
  \node[below] at (O) {$O$};
  \node[below] at (D) {$D$};
  \node[below] at (E) {$E$};
\end{tikzpicture}
\end{center}
\end{problem}
\clearpage


\begin{problem}{2.1.10}
质量为 $m$ 的小环套在绳上, 绳子跨过定滑轮并在另一端系一质量为 $M$ 的物体. 若绳子不可伸长、不计质量, 绳子与滑轮间的摩擦力可忽略, 小环相对于绳子以加速度 $a$ 下落, 求小环与绳子间的摩擦力.
\begin{center}
\begin{tikzpicture}[>=Stealth, thick]
  % Left FBD: block M
  \draw (-0.42,-0.30) rectangle (0.42,0.30);
  \node[left=2pt] at (-0.42,0) {$M$};
  \draw[->] (0,0.30) -- (0,1.25) node[right=1pt] {$T$};
  \draw[->] (0,-0.30) -- (0,-1.25) node[below=1pt] {$Mg$};
  \draw[->] (0.95,0.55) -- (0.95,-0.75) node[below=1pt] {$a_M$};

  % Right FBD: ring m
  \draw (2.8,0) ellipse (0.38 and 0.24);
  \node[right=2pt] at (3.18,0) {$m$};
  \draw[->] (2.8,0.24) -- (2.8,1.25) node[right=1pt] {$f$};
  \draw[->] (2.8,-0.24) -- (2.8,-1.25) node[below=1pt] {$mg$};
  \draw[->] (3.85,0.55) -- (3.85,-0.75) node[below=1pt] {$a - a_M$};
\end{tikzpicture}

\end{center}
\end{problem}
\clearpage


\begin{problem}{2.1.11}
一个质量为 $m$ 的物体通过一根质量可以不计的绳子绕水平棒 $1\frac{1}{4}$ 周后于另一端加一水平力 $F$ , 如图2.10所示. 若绳子和棒之间的摩擦因数为 $\mu$ , 要使物体保持静止状态, 应施加多大的水平拉力?
\begin{center}
\begin{tikzpicture}[>=Stealth, thick]
  % Hanging block m
  \draw (-1.00,-2.35) rectangle (-0.08,-1.65);
  \node at (-0.54,-2.00) {$m$};

  % Vertical chain links from block up to rod
  \foreach \y in {-1.55,-1.35,-1.15,-0.95,-0.75,-0.55,-0.35,-0.16} {
    \draw[fill=white] (-0.54,\y) ellipse (0.065 and 0.105);
  }

  % Back passes of the chain behind the rod
  \draw[fill=white] (-0.44,-0.02) ellipse (0.07 and 0.10);
  \draw[fill=white] (-0.14,0.10) ellipse (0.07 and 0.10);

  % Tilted horizontal rod
  \begin{scope}[rotate=22]
    \fill[white] (-2.1,-0.14) rectangle (1.9,0.14);
    \draw (-2.1,0.14) -- (1.9,0.14);
    \draw (-2.1,-0.14) -- (1.9,-0.14);
  \end{scope}

  % Front passes wrapping around the rod (1.25 turns)
  \draw[fill=white, rotate around={-35:(-0.46,-0.10)}] (-0.46,-0.10) ellipse (0.065 and 0.10);
  \draw[fill=white, rotate around={-35:(-0.32,0.04)}] (-0.32,0.04) ellipse (0.065 and 0.10);
  \draw[fill=white, rotate around={-35:(-0.19,-0.12)}] (-0.19,-0.12) ellipse (0.065 and 0.10);
  \draw[fill=white, rotate around={-35:(-0.06,0.03)}] (-0.06,0.03) ellipse (0.065 and 0.10);
  \draw[fill=white] (-0.54,0.11) circle (0.07);
  \draw[fill=white] (-0.23,0.18) circle (0.07);

  % Rightward chain and pulling force F
  \foreach \i in {0,...,9} {
    \pgfmathsetmacro{\cx}{0.10 + \i*0.20}
    \pgfmathsetmacro{\cy}{0.04 - \i*0.055}
    \draw[fill=white, rotate around={-15.5:(\cx,\cy)}] (\cx,\cy) ellipse (0.105 and 0.055);
  }
  \draw[->] (2.08,-0.51) -- (2.80,-0.71) node[below=1pt] {$F$};
\end{tikzpicture}

\end{center}
\end{problem}
\clearpage


\begin{problem}{2.1.12}
一根质量为 $M$ 的均质绳子, 两端固定在同一高度的两个钉子上, 在其中点挂一质量为 $m$ 的物体. 设 $\alpha, \beta$ 分别为在绳子中点和端点处绳子的切线方向与竖直方向的夹角, 如图 2.13 所示. 证明:

\[
\frac {\tan \alpha}{\tan \beta} = \frac {M + m}{m}
\]
\begin{center}
\begin{tikzpicture}[>=Stealth, thick]
  % Rope curve from center cusp (0,0) to both nails and tangent extensions
  \draw (0,0) .. controls (-0.85,0.42) and (-1.28,0.95) .. (-1.6,1.8) -- (-1.97,2.78);
  \draw (0,0) .. controls (0.85,0.42) and (1.28,0.95) .. (1.6,1.8) -- (1.97,2.78);

  % Nails
  \fill (-1.6,1.8) circle (2.2pt) node[right=2pt] {钉子};
  \fill (1.6,1.8) circle (2.2pt) node[left=2pt] {钉子};

  % Vertical dashed reference lines and beta angles at nails
  \draw[dashed] (-1.6,1.8) -- (-1.6,2.80);
  \draw[dashed] (1.6,1.8) -- (1.6,2.80);
  \node at (-1.77,2.64) {$\beta$};
  \node at (1.77,2.64) {$\beta$};

  % Vertical dashed reference and tangents at bottom cusp (alpha angles)
  \draw[dashed] (0,0) -- (0,0.85);
  \draw (0,0) -- (-0.95,0.47);
  \draw (0,0) -- (0.95,0.47);
  \node at (-0.28,0.36) {$\alpha$};
  \node at (0.28,0.36) {$\alpha$};

  % Hanging mass m
  \draw (0,0) -- (0,-0.92);
  \draw (-0.52,-0.92) rectangle (0.52,-1.52);
  \node at (0,-1.22) {$m$};
\end{tikzpicture}

\end{center}
\end{problem}
\clearpage


\begin{problem}{2.1.13}
一质量为 $M$ 的均质链条套在一表面光滑、顶角为 $\alpha$ 、底面处于水平的正圆锥上保持静止，若链条平面平行于水平面，求链条中的张力.
\end{problem}
\clearpage


\begin{problem}{2.1.14}
真空中有缺陷的、质量为 950kg 的卫星，被飞船用一根 50m 长、线密度为 1kg/m 的均质绳牵着以 $5m/s^{2}$ 的加速度做匀加速直线运动：

(1) 飞船作用在绳上的力多大?

(2) 计算绳的张力；

(3) 飞船上的人精疲力竭睡着了,飞船助推器的控制电路之一发生短路,使加速度变为 $1 \, m/s^{2}$ 的减速度.这个事故会发生什么后果?
\end{problem}
\clearpage


\begin{problem}{2.1.15}
某人观察到一空间站始终停留在地球同一点的垂直上空. 问观察者位于地球的什么地方? 尽可能详细地描述一下此空间站的轨道.
\end{problem}
\clearpage


\begin{problem}{2.1.16}
公园里有一种转动水平圆盘, 孩子可以坐在盘面上任何地方. 当盘开始加快转速时, 如果摩擦力不够大, 孩子就可能滑出去. 设孩子的质量为 50kg, 摩擦因数为 0.4, 盘的角速度为 2rad/s, 问孩子坐在盘上不被滑开的最大半径多大?
\end{problem}
\clearpage


\begin{problem}{2.1.17}
一小物体放在一半径为 $R$ 的水平圆盘上, 小物体与圆盘间的静摩擦因数为 $\mu$ . 若圆盘绕其轴的角速度逐渐增大到一个值时, 小物体滑出圆盘并最终落到比盘面低 $h$ 的地面上. 问从它离开圆盘的那一点算起, 小物体越过的水平距离多大?
\end{problem}
\clearpage


\begin{problem}{2.1.18}
太阳离银河系中心大约相距 $2.5 \times 10^{4}$ 光年, 近似地以 $1.7 \times 10^{8}$ 年的周期绕此中心做圆周运动. 地球离太阳的距离为 8 光分. 用以上数据, 求出以太阳质量作为单位的银河系的近似质量. 假定作用在太阳上的引力近似为银河系的质量集中在其中心时对太阳的引力.
\end{problem}
\clearpage


\begin{problem}{2.1.19}
用以下近似数据计算地球和太阳的平均密度之比：

$\theta$ 从地球上看太阳的角直径, $\theta=0.5^{\circ}$ ;

l 为地球表面上纬度为 $1^{\circ}$ 的长度, l=100km;

t 为地球公转周期, t=1 年=$3\times 10^{7}$s:

g 为重力加速度, $g=10ms^{-2}$
\end{problem}
\clearpage


\begin{problem}{2.1.20}
(1) 一个均质的球形物体以角速度 $\omega$ 绕对称轴转动, 如果仅仅靠自身的万有引力阻碍球体的离心分解, 该物体必须具有的最小密度多大? 用这一点估计巨蟹座中转速为每秒 30 转的脉冲星的最小密度. (这是公元 1054 年在中国广泛地被观察到的一个超新星爆发的遗迹!)

(2) 如果脉冲星的质量与太阳的质量相当, 约为 $2 \times 10^{30} ~kg$ , 这个脉冲星的最大可能半径多大?

(3) 若脉冲星的密度与核物质的密度相当, 它的半径多大?
\end{problem}
\clearpage


\begin{problem}{2.1.21}
一顶角为 $2\alpha (\alpha > 45^{\circ})$ 的倒立的圆锥面，以恒定角速度 $\omega$ 绕其对称轴旋转，在内表面上距轴为 $r$ 处有一质点，若质点与锥面间的摩擦因数为 $\mu$ ，要使质点相对锥面静止， $\omega$ 应处于什么范围内？
\end{problem}
\clearpage


\begin{problem}{2.1.22}
一圆柱形刚性杆上套一质量为 $m$ 的小环，杆绕通过一端的固定竖直轴以恒定角速度 $\omega$ 旋转，旋转时杆与竖直轴的夹角为 $\alpha$ 。若小环与杆间的摩擦因数为 $\mu$ ，以小环距杆的固定端的距离 $x$ 为坐标列出小环的运动微分方程。
\end{problem}
\clearpage


\begin{problem}{2.1.23}
质量为 m 的小球通过一根原长为 l、劲度系数为 k 的弹性绳系在光滑的水平台面上的 P 点，水平台面绕离 P 点距离为 b 的固定竖直轴以恒定的角速度 $\omega$ 转动。列出小球在弹性绳拉直状态下的运动微分方程。
\end{problem}
\clearpage


\begin{problem}{2.1.24}
质量为 M 的斜劈放在摩擦因数为 $\mu$ 的粗糙水平面上, 质量为 $m_{1}$ 的物体用不计质量、不可伸长的弦线挂起并跨过固连于斜劈的光滑滑轮与在无摩擦的斜劈面上滑动的质量为 $m_{2}$ 的物体相连, 斜劈面的倾角为 $\theta$ . 如图 2.17 所示.

(1) 求当 $\mu$ 非常大时 $m_{1}$ 和 $m_{2}$ 的加速度以及弦线的张力；

(2) 求使斜劈能保持静止的最小摩擦因数；

(3) 若 $\mu$ 不够大, 斜劈不能保持静止, 有向左的加速度 $a_{M}$ , 求此 $a_{M}$ , 列出足够的方程即可, 不必解出 $a_{M}$ . 设滑轮很小, $m_{1}, M$ 紧靠时, 连结 $m_{1}$ 的弦线仍可视为在竖直方向.
\begin{center}
\begin{tikzpicture}[>=Stealth, thick]
  % Rough horizontal ground
  \draw (-1.1,0) -- (4.1,0);
  \foreach \x in {-0.9,-0.6,...,3.9} {
    \draw[thin] (\x,0) -- (\x-0.25,-0.25);
  }
  \node at (1.5,-0.48) {$\mu$};

  % Wedge M
  \draw (0,0) -- (3.6,0) -- (0,2.0) -- cycle;
  \node at (1.15,0.68) {$M$};
  \node at (2.82,0.22) {$\theta$};

  % Pulley at top-left corner
  \draw (0,2.0) -- (-0.14,2.14);
  \draw[fill=white] (-0.14,2.14) circle (0.22);

  % Hanging mass m1
  \draw (-0.36,2.14) -- (-0.36,1.08);
  \draw (-0.58,0.66) rectangle (-0.14,1.08);
  \node[left=2pt] at (-0.58,0.87) {$m_1$};

  % Block m2 on the incline and string
  \begin{scope}[shift={(1.512,1.16)}, rotate=-29.05]
    \draw (0,0) rectangle (0.56,0.52);
    \node at (0.28,0.78) {$m_2$};
    \coordinate (m2attach) at (0,0.26);
  \end{scope}
  \draw (-0.033,2.332) -- (m2attach);
\end{tikzpicture}

\end{center}
\end{problem}
\clearpage


\begin{problem}{2.1.25}
一质量为 M 的、上表面水平的楔形物体，在倾角为 $\alpha$ 的固定的光滑斜面上运动，楔形物体上放一质量为 m 的质点，如图 2.20 所示。若 m、M 间无摩擦，求：

(1) $m$ 相对于 $M$ 的加速度；

(2) $M$ 对斜面的压力.
\begin{center}
\begin{tikzpicture}[>=Stealth, thick]
  % Horizontal ground
  \draw (-1.0,0) -- (4.5,0);
  \foreach \x in {-0.8,-0.5,...,4.3} {
    \draw[thin] (\x,0) -- (\x-0.25,-0.25);
  }

  % Fixed inclined plane
  \draw (-0.8,2.4) -- (4.3,0);
  \node at (3.35,0.20) {$\alpha$};

  % Wedge M with horizontal top surface
  \draw (0.369,1.85) -- (2.6,1.85) -- (2.6,0.80);
  \node at (1.98,1.45) {$M$};

  % Block m on top of wedge M
  \draw (1.15,1.85) rectangle (1.95,2.45);
  \node at (1.55,2.15) {$m$};
\end{tikzpicture}

\end{center}
\end{problem}
\clearpage


\begin{problem}{2.1.26}
图 2.22 中水平桌面上的滑块质量为 M，一条其两端分别系着质量为 $m_{1}$ 、 $m_{2}$ 

的两个物体的不计质量、不可伸长的细绳套在滑块上，滑块、绳子与桌面相互间均无摩擦。试证明滑块的加速度为

\[
a = \frac {4 m _ {1} m _ {2}}{M (m _ {1} + m _ {2}) + 4 m _ {1} m _ {2}} g
\]
\begin{center}
\begin{tikzpicture}[>=Stealth, thick]
  % Table legs and apron (drawn first so hanging strings pass cleanly in front)
  \draw (-2.05,0.16) -- (-2.05,-2.05) -- (-1.85,-2.10) -- (-1.85,0.11);
  \draw (-1.85,-2.10) -- (-1.73,-2.00) -- (-1.73,-0.55);
  \draw (1.02,-0.546) -- (1.02,-2.45) -- (1.22,-2.50) -- (1.22,-0.592);
  \draw (1.22,-2.50) -- (1.34,-2.40) -- (1.34,-0.51);

  % Apron under tabletop
  \draw (-1.85,-0.52) -- (1.02,-1.18);
  \draw (1.34,-1.08) -- (2.45,-0.22);

  % Tabletop surface and front/side thickness
  \draw (-0.85,1.40) -- (-2.20,0.35) -- (1.30,-0.45) -- (2.65,0.60);
  \draw (-2.20,0.35) -- (-2.20,0.19) -- (1.30,-0.61) -- (2.65,0.44);
  \draw (1.30,-0.45) -- (1.30,-0.61);

  % Slider M (ellipse on tabletop) and string wrapped around it
  \draw[rotate around={-13:(0.47,0.954)}] (0.47,0.954) ellipse (0.503 and 0.25);
  \node at (1.18,1.32) {$M$};
  \draw (-0.02,1.066) -- (-1.22,0.126) -- (-1.22,-1.25);
  \draw (0.96,0.842) -- (-0.24,-0.098) -- (-0.24,-1.58);

  % Hanging masses m1 and m2
  \draw[fill=white] (-1.42,-1.25) rectangle (-1.02,-1.58);
  \node[right=2pt] at (-1.02,-1.42) {$m_1$};
  \draw[fill=white] (-0.44,-1.58) rectangle (-0.04,-1.91);
  \node[right=2pt] at (-0.04,-1.75) {$m_2$};
\end{tikzpicture}

\end{center}
\end{problem}
\clearpage


\begin{problem}{2.1.27}
$q_{1}, q_{2}$ 轴在竖直的 $xy$ 平面内，与 $x$ 轴的夹角分别为 $\alpha$ 和 $\beta$ . 一个抛射体在 $xy$ 平面内运动，今取 $q_{1}, q_{2}$ 为坐标，给出或导出抛射体的运动微分方程. 重力沿 $y$ 轴的负方向.
\end{problem}
\clearpage


\begin{problem}{2.1.28}
质量分别为 $M$ 和 $M + m$ 的两个人, 分别拉住跨过定滑轮两边的绳子往上爬. 开始时两人与滑轮的距离都是 $h$ , 都从静止开始爬. 设绳子不可伸长, 绳子质量与绳子、滑轮间的摩擦力均可不计. 证明: 如质量较轻的人在 $t$ 秒内爬到滑轮, 这时质量较重的人与滑轮的距离为 $\frac{m}{M + m}\left(h + \frac{1}{2} gt^2\right)$ .
\end{problem}
\clearpage


\begin{problem}{2.1.29}
一质量为 m 的跳水运动员, 初速为零地从 10m 跳台上跳下.

(1) 求入水速度和从起跳到入水大致所用的时间；

(2) 假定作用在跳水者身上水的浮力和他所受的重力正好抵消, 作用在他身上的水的黏滞力大小为 $bv^2$ , 试列出跳水员在水中垂直下沉的运动微分方程, 以边界条件 $x = 0$ 处 $v = v_0$ , 求解速度 $v$ 作为水面下深度 $x$ 的函数;

(3) 若 $\frac{b}{m}=0.4m^{-1}$ ，求当 $v=\frac{1}{10}v$ 时的深度；

(4) 求解跳水员在水下的垂直深度作为在水下的时间的函数.
\end{problem}
\clearpage


\begin{problem}{2.1.30}
一个质量为 $m$ 的质点以初速 $v$ 在 $y$ 方向投射到沿 $x$ 方向以不变速度 $V$ 运动着的水平皮带上并留在上面运动. 若质点与皮带间的滑动摩擦因数为 $\mu$ , 质点接触皮带的初位置取为固定的 $xy$ 坐标系的原点, 求质点在皮带上刚停止滑动时的 $xy$ 坐标.
\end{problem}
\clearpage


\begin{problem}{2.1.31}
在水平的 $xy$ 平面内有一留声机转盘, 它以恒定的角速度 $\omega$ 绕过盘心的铅垂轴旋转. 在转盘上滑动的物体在水平方向受到两个真实力作用: 一个是大小为 $kr(r$ 为质点至盘心的距离, $k$ 为常量), 方向指向盘心的弹性力; 另一个是摩擦力, 与相对速度成正比, 比例系数为 $c, c$ 是正值常量, 物体可视为质点.

(1) 如物体待在盘上任何位置均能相对转盘保持静止, 问 k 有多大?

(2) 假设 k 是(1)问所得之值, 在一般的初始条件下, 求物体的相对速度 $\dot{x}$ 、 $\dot{y}$ ;

(3) 用(2)的结果进而求 $x(t)$ 和 $y(t)$ .
\end{problem}
\clearpage


\begin{problem}{2.1.32}
一个质点垂直上抛, 初速为 $v_{0}$ , 如空气阻力与速率的平方成正比, 证明: 回到初始位置时的速率为 $\frac{v_{0} v_{t}}{(v_{0}^{2} + v_{t}^{2})^{1/2}}$ , 其中 $v_{t}$ 是极限速率.
\end{problem}
\clearpage


\begin{problem}{2.1.33}
两个质量均为 $m$ 的质点 $A$ 和 $B$ 连在一个劲度系数为 $k$ 的弹簧的两端. 开始两质点静放在光滑的水平面上, 弹簧处于原长, 然后沿 $AB$ 方向给 $B$ 以恒力 $ka$ . 求两质点的运动学方程.
\end{problem}
\clearpage


\begin{problem}{2.1.34}
质量为 $m$ 的质点在阻力等于 $mk(v^3 + a^2 v)$ 的介质中运动, 其中 $k$ 和 $a$ 均为常量, 初速度为 $v_0$ . 求其路程与时间的关系, 并由此导出 $k \to 0$ 时的关系来检验所得结果的正确性.
\end{problem}
\clearpage


\begin{problem}{2.1.35}
质量为 $m$ 的质点静止地连在一弹簧的端点，弹簧的劲度系数为 $k$ ，并悬挂在固定支点上。当 $t = 0$ 时，对质点施一向下的恒力 $F$ ，作用时间为 $t_0$ 。证明：当外力去掉后，质点由其平衡位置 $x = x_0$ 所做的位移为

\[
x - x _ {0} = \frac {F}{k} [ \cos \omega (t - t _ {0}) - \cos \omega t ]
\]

其中 $\omega^{2}=\frac{k}{m}.$
\end{problem}
\clearpage


\begin{problem}{2.1.36}
空间内有均匀的电场和磁场, $E=Ei,B=B(\cos\alpha i+\sin\alpha j)$ ,$\alpha$为常量.具有质量为m、电量为q的粒子在坐标系原点以初速 $v_{0}=\dot{x}_{0}i+\dot{y}_{0}j+\dot{z}_{0}k$ 开始运动.若重力可忽略不计,求质点以后的运动.
\end{problem}
\clearpage


\begin{problem}{2.1.37}
具有质量 m 和电量 q 的质点, 以速度 $V_{0}$ 垂直于电场方向进入按 E = $E_{0}\cos\omega t$ 规律变化的均匀电场, 若质点还受到一个跟速度一次幂成比例的阻力 R = -$\gamma$v, $\gamma$ 为正值常量. 求质点的运动(不计重力的作用).
\end{problem}
\clearpage


\begin{problem}{2.1.38}
一长度为 L 的水平封闭圆管绕其一端的竖直轴以恒定的角速度 $\omega$ 旋转，管内装有密度为 $\rho_{0}$ 的液体及一截面半径略小于管内径的长为 l、密度为 $\rho(\rho > \rho_{0})$ 的小圆柱体。开始时，小圆柱体紧靠轴并相对于圆管静止。问经多长时间小圆柱体到达管的另一端。
\end{problem}
\clearpage


\begin{problem}{2.1.39}
一个质点以初始速率 $v_{0}$ 、仰角 $\alpha$ 抛出，若空气阻力与速率的平方成正比，比例系数为 $mk, m$ 为质点的质量。证明：

(1) 运动微分方程可以写成

\[
\ddot {x} = - k \dot {x} \dot {s}, \quad \ddot {y} = - k \dot {y} \dot {s} - g
\]

其中 x 轴沿水平方向，y 轴竖直向上， $\dot{s}^{2}=\dot{x}^{2}+\dot{y}^{2}$ ;

(2) 轨道微分方程可写为

\[
\frac {\mathrm{d} ^ {2} y}{\mathrm{d} x ^ {2}} = - g \mathrm{e} ^ {2 k s} / (v _ {0} \cos \alpha) ^ {2}
\]

其中 s 是自抛射点沿轨道经过的路程；

(3) 忽略空气阻力时,积分上述轨道微分方程,可得抛物线形式的轨道.
\end{problem}
\clearpage


\begin{problem}{2.1.40}
有一光滑旋转抛物面,其方程为

\[
z = - \frac {1}{2} (x ^ {2} + y ^ {2})
\]

z 轴竖直向上. 在其顶点有一质量为 m 的小物体, 受到微小扰动后自静止开始下滑. 求抛

物面对物体的作用力(取质点运动平面为 $xz$ 平面.)
\end{problem}
\clearpage


\begin{problem}{2.1.41}
一个 $3:4:5$ 的斜面固连在一转盘上, 如图 2.28 所示, 一木块静止在斜面上, 斜面和木块之间的静摩擦因数 $\mu_{s} = \frac{1}{4}$ . 求此木块能保持在离转盘中心的水平距离为 $40\mathrm{cm}$ 处相对转盘不动的最小转动角速度 $\omega$ .
\begin{center}
\begin{tikzpicture}[>=Stealth, thick]
  % Lower shaft
  \draw (0,-1.07) -- (0,-1.92);

  % Turntable disk (3D ellipse + rim)
  \draw (-2.6,0) -- (-2.6,-0.22) arc (180:360:2.6 and 0.85) -- (2.6,0);
  \draw[fill=white] (0,0) ellipse (2.6 and 0.85);

  % Central upper shaft
  \draw (0,0) -- (0,1.88);

  % Rotation arrow omega
  \draw[->, very thick] (-0.62,-0.46) arc (238:302:1.2 and 0.38);
  \node at (0.05,-0.24) {$\omega$};

  % 3:4:5 wedge on turntable
  \coordinate (W1) at (-2.10,-0.05);
  \coordinate (W2) at (-0.95,-0.05);
  \coordinate (W3) at (-2.10,0.81);
  \coordinate (W2b) at (-0.50,0.35);
  \coordinate (W3b) at (-1.65,1.21);

  \fill[white] (W1) -- (W2) -- (W2b) -- (W3b) -- (W3) -- cycle;
  \draw (W1) -- (W2) -- (W3) -- cycle;
  \draw (W3) -- (W3b) -- (W2b) -- (W2);

  \node[left=3pt] at (-2.12,0.52) {$3$};
  \node[below=2pt] at (-1.52,-0.05) {$4$};
  \node at (-1.24,0.40) {$5$};

  % Small block resting on the sloped face
  \filldraw[fill=white] (-1.62,0.72) -- (-1.45,0.62) -- (-1.34,0.72) -- (-1.51,0.82) -- cycle;
  \filldraw[fill=white] (-1.62,0.72) -- (-1.62,0.86) -- (-1.45,0.76) -- (-1.45,0.62) -- cycle;
  \filldraw[fill=white] (-1.45,0.62) -- (-1.45,0.76) -- (-1.34,0.86) -- (-1.34,0.72) -- cycle;
  \filldraw[fill=white] (-1.62,0.86) -- (-1.45,0.76) -- (-1.34,0.86) -- (-1.51,0.96) -- cycle;

  % 40 cm horizontal distance indicator
  \draw (-1.50,0.98) -- (-1.50,1.88);
  \draw[<->] (-1.50,1.68) -- (0,1.68) node[midway, above=1pt] {$40\,\mathrm{cm}$};
\end{tikzpicture}

\end{center}
\end{problem}
\clearpage


\begin{problem}{2.1.42}
一水平转盘由静止开始启动，并以匀角加速度 $0.04\pi \mathrm{rad} / \mathrm{s}^2$ 加速。一小孩坐在离转盘 $6\mathrm{m}$ 远的椅子上，手中握着一个 $2\mathrm{kg}$ 的球。求转盘启动5s后那一瞬间，小孩为握住球必须施以的力的大小和方向。
\end{problem}
\clearpage


\begin{problem}{2.1.43}
一个在竖直平面内半径为 $a$ 的圆环，以角速度 $\omega$ 绕其竖直直径旋转。一质量为 $m$ 的小球在环上无摩擦地滑动，如图2.30所示。

(1) 在什么特殊条件下, 小球在 $\theta=0$ 处是稳定平衡?

(2) 当 $\omega$ 为何值时, 小球可处于另一稳定平衡位置, 求出这时的 $\theta$ 值.
\begin{center}
\begin{tikzpicture}[>=Stealth, thick]
  % Circular hoop
  \draw (0,0) circle (1.6);

  % Vertical rotation axis
  \draw (0,-1.6) -- (0,2.35);

  % Rotation arrow omega
  \draw[->] (-0.28,2.08) arc (140:395:0.42 and 0.14);
  \node[right=2pt] at (0.38,2.14) {$\omega$};

  % Radius to bead m at angle theta = 40 deg from downward vertical
  \draw (0,0) -- (-50:1.6);
  \fill (-50:1.6) circle (3pt) node[below right=1pt] {$m$};
  \node at (0.30,-0.68) {$\theta$};
\end{tikzpicture}

\end{center}
\end{problem}
\clearpage


\begin{problem}{2.1.44}
一光滑的水平圆盘以角速度 $\omega$ 绕通过盘心的垂直轴转动. 在圆盘上位于离盘心距离 $R$ 的人给质量为 $m$ 的光滑硬币（忽略大小）朝向盘心方向的推力，使硬币得到相对于圆盘的初速度 $v$ . 证明：在圆盘上的人看来，在一段时间内 $[( \omega t)^2$ 可以忽略]运动是一条抛物线，给出这条抛物线的方程.
\end{problem}
\clearpage


\begin{problem}{2.1.45}
卫星绕地球做圆轨道运动,卫星内的宇航员将一个小物体放在卫星质心和地心连线上离卫星质心 $\Delta r$ 处然后释放.求在卫星参考系中宇航员所看到的释放后物体的运动(设卫星像月球那样朝向地球的部分始终朝向地球).
\end{problem}
\clearpage


\begin{problem}{2.1.46}
一质量为 $m$ 的质点用原长为 $l$ 、劲度系数为 $k$ 的弹簧挂起来， $t = 0$ 时，弹簧上端的支撑点开始上下按正弦函数做振动，振动的振幅为 $A$ ，角频率为 $\omega (\omega \neq \sqrt{k / m})$ 。求质点的运动。
\end{problem}
\clearpage


\begin{problem}{2.1.47}
一块光滑的水平板绕过板上一点的竖直轴做恒定角速度为 $\omega$ 的转动. 一质点用一根原长为零、劲度系数为 $k$ 的弹簧连接于轴与平板的交点. 若取开始时质点与交点 (取为坐标系原点) 的连线为 $x$ 轴, $z$ 轴沿 $\omega$ 的方向, 质点开始在 $(a,0)$ 处以 $(0,v_0)$ 的速度运动, 求质点的运动学方程 $x(t), y(t)$ .
\end{problem}
\clearpage


\begin{problem}{2.1.48}
考虑一组荷质比 e/m 都相同的带电粒子, 证明: 这些粒子在一个小的磁场 B (常矢量) 中的运动同在一个适当选择的角速度 $\omega$ 转动的参考系中看到的没有这个磁场的运动相同, 求这相应的 $\omega$ 值, 并对“小场”作出说明.
\end{problem}
\clearpage


\begin{problem}{2.1.49}
一个质量为 $m$ 的质点在光滑的水平面上以角速度 $\omega$ 做匀速率圆周运动, 其向心力是通过一劲度系数为 $k$ 的弹簧提供的. 使质点突然获得一个很小的径向速度分量, 求所引起的径向振动的频率.
\end{problem}
\clearpage


\begin{problem}{2.1.50}
一个质量为 $M$ 的小珠套在光滑的水平的金属丝上, 另一个质量为 $m$ 的质点通过无质量的不可伸长的、长度为 $a$ 的绳子连在小珠上. 握住小珠与质点, 使绳子沿金属丝拉直, 然后突然释放. 求此后绳子张力与绳子、竖直线间夹角 $\theta$ 的关系.
\end{problem}
\clearpage


\begin{problem}{2.1.51}
质量分别为 $m$ 和 $M$ 的两个小环可在一光滑的竖直平面内的圆圈上运动，两环被一根无质量的、不可伸长的绳子连着。证明：只要绳子保持拉紧，绳子张力为 $2mMg\tan \alpha \cos \theta /(m + M)$ ，其中 $2\alpha$ 是绳子拉紧时对圈心所张的角， $\theta$ 是绳子与水平线间的夹角， $g$ 为重力加速度。
\end{problem}
\clearpage


\begin{problem}{2.1.52}
三个点源等距离地固定在一个半径为 $a$ 、圆心为原点的圆周上，每个点源对质量为 $m$ 的质点的作用力为引力 $F = -kR$ ，其中 $k$ 为正值常量， $R$ 为质点对点源的位矢。质点在 $t = 0$ 时放入力场，初始条件为 $r = r_0 \cdot r = v_0$ ，如图2.34所示。

(1) 求解质点的运动, 用 $r_0, v_0$ 及其他参数表示 $\boldsymbol{r}(t)$ ;

(2) 在什么条件下, 质点的运动轨道为圆?
\begin{center}
\begin{tikzpicture}[>=Stealth, thick]
  % Dashed circle of radius a
  \draw[dashed] (0,0) circle (1.85);
  \node[below left=1pt] at (0,0) {$O$};

  % Three equally spaced point sources and dashed radii
  \foreach \ang in {30,150,-90} {
    \draw[dashed] (0,0) -- (\ang:1.85);
    \fill (\ang:1.85) circle (2.8pt);
  }
  \node at (-0.95,0.34) {$a$};
  \node at (0.95,0.34) {$a$};
  \node at (0.18,-0.95) {$a$};

  % Particle m, position vector r0, and initial velocity v0
  \coordinate (M) at (-0.52,0.98);
  \draw[->] (0,0) -- (M) node[midway, right=1pt] {$\boldsymbol{r}_0$};
  \node[left=2pt] at (M) {$m$};
  \draw[->] (M) -- (-0.08,1.50) node[right=1pt] {$\boldsymbol{v}_0$};
\end{tikzpicture}

\end{center}
\end{problem}
\clearpage


\begin{problem}{2.1.53}
$n$ 个质点质量均为 $m$ , 任何两个质点间相互作用力为斥力, 大小与它们间的距离成正比, 比例系数为 $\mu m^2$ . 开始把它们静止地放在一个光滑的水平面上处于以 $O$ 为中心、 $a$ 为半径的圆圈上的对称位置. 证明: 同时释放后, 任一质点位于离 $O$ 点距离 $r$ 处时的速率为

$\left[\mu nm(r^2 - a^2)\right]^{1/2}$ 提示： $2\left[\sin^2 \frac{\pi}{n} + \sin^2 \frac{2\pi}{n} + \cdots + \sin^2 \frac{(n-1)\pi}{n}\right] = n$ （ $n \geqslant 2$ ).
\end{problem}
\clearpage


\begin{problem}{2.1.54}
考虑一质量为 $m$ 的粒子在力 $\bm{F} = -kr$ 作用下运动，这里 $k$ 是正值常量， $\bm{r}$ 是粒子的位矢.

(1) 证明粒子在一平面内运动；

(2) 设在 $t = 0$ 时, $x = a, y = 0, \dot{x} = 0, \dot{y} = v_0$ , 求 $x(t)$ 和 $y(t)$ ;

(3) 证明该轨道是一椭圆；

(4) 求出周期.
\end{problem}
\clearpage


\begin{problem}{2.1.55}
质量分别为 $m$ 及 $m'$ 的两个质点，用未伸缩时长度为 $a$ 的弹性绳相连，绳的劲度系数 $k = \frac{2mm' \omega^2}{m + m'}$ . 将此系统放在光滑的水平管内，管子绕通过管上某点的竖直轴以恒定角速度 $\omega$ 转动. 开始时，两质点相对于管子是静止的，两质点间距为 $a$ . 试求此后任何时刻两质点间的距离.
\end{problem}
\clearpage


\begin{problem}{2.1.56}
质量为 $m$ 的小环套在半径为 $a$ 的光滑圆圈上，可在其上滑动。如圆圈在水平面内以恒定角速度 $\omega$ 绕圈上某点 $O$ 转动，转轴垂直于水平面。试求小环沿圆圈切线方向的运动微分方程。
\end{problem}
\clearpage


\begin{problem}{2.1.57}
两个粒子质量为 $m$ 、带正电荷 $q$ ，用劲度系数为 $k$ 、平衡长度为零的弹簧相连。沿水平方向有恒定电场 $E = E_0 i$ ，考虑粒子之间的库仑力，忽略磁效应、辐射效应和相对论效应等，假定粒子没有碰撞。

(1) 如果粒子在水平 x 方向的光滑直线上滑动, 它们之间的距离 d 不变. 求 d;

(2) 若(1)中两粒子间的距离 $d(t)$ 为在一个平衡位置附近的小振动, 求频率;

(3) 若粒子在水平的光滑桌面上滑动, 求运动微分方程的通解, 可在答案中保留积分.
\end{problem}
\clearpage


\begin{problem}{2.1.58}
一个半径为 $R$ 的半球形碗以不变角速度 $\Omega$ 绕竖直的对称轴转动. 质量为 $M$ 的一个粒子在重力影响下在碗的内表面运动, 粒子除受到碗面的支持力外, 还受到一个摩擦力 $F = -kv$ , 其中 $k$ 为正值常量, $v$ 是相对于碗的速度.

（1）用球坐标 $\theta, \varphi$ （坐标原点取在碗心），写出粒子相对于转动参考系的运动微分方程. 证明：如果 $\Omega > \sqrt{\frac{g}{R}}$ ，除了碗底这个平衡位置以外，存在第二个平衡位置，并求出这第二个平衡位置的 $\theta$ 值.

(2) 考虑粒子在碗底这个平衡位置附近的运动。在此平衡点取一个固定的笛卡儿坐标系 $x, y, z$ ，除计算重力提供的恢复力以外，不考虑碗的曲率。证明：对于 $|x| \ll R$ 、 $|y| \ll R$ ，粒子的一个特解为

\[
x = \operatorname{Re} (x _ {0} \mathrm{e} ^ {\lambda t}), \quad y = \operatorname{Re} (y _ {0} \mathrm{e} ^ {\lambda t})
\]

其中 $\lambda$ 满足

\[
\left(\lambda^ {2} + \frac {k}{M} \lambda + \frac {g}{R}\right) ^ {2} + \left(\frac {k}{M}\right) ^ {2} \Omega^ {2} = 0
\]

(3) 求碗的角速度 $\Omega_{0}$ ，对于此角速度，粒子的运动是周期性的.
\end{problem}
\clearpage


\begin{problem}{2.1.59}
一个质点在重力作用下沿一粗糙的竖直的圆周运动. 如从水平直径的一端的静止状态开始运动, 到圆的最低点时刚好静止. 求摩擦因数 $\mu$ 必须满足的关系.
\end{problem}
\clearpage


\begin{problem}{2.1.60}
一个质量为 $M$ 的小球, 用一根劲度系数为 $k$ 、自然长度为零的无质量弹簧悬挂于天花板下, 小球下面挂一根相同的弹簧, 原处于平衡状态. 今有一个力 $F(t)$ 作用于下面弹簧的下端, 如图 2.42 所示. 求上面弹簧的长度 $x_{1}$ 与 $F(t)$ 的关系; 如 $F(t) = \lambda t$ , $\lambda$ 为常量, 求 $x_{1}(t)$ .
\begin{center}
\begin{tikzpicture}[>=Stealth, thick]
  % Ceiling with hatching
  \draw (-1.35,0) -- (1.45,0);
  \foreach \x in {-1.25,-1.00,...,1.25} {
    \draw[thin] (\x,0) -- (\x+0.22,0.22);
  }

  % Upper spring k
  \draw (0,0) -- (0,-0.28);
  \draw[decorate, decoration={zigzag, segment length=5.5pt, amplitude=3.5pt}] (0,-0.28) -- (0,-1.24);
  \node[left=6pt] at (0,-0.76) {$k$};
  \draw (0,-1.24) -- (0,-1.52);

  % Ball M
  \draw (0,-1.90) circle (0.38);
  \node at (0,-1.90) {$M$};

  % Lower spring k
  \draw (0,-2.28) -- (0,-2.56);
  \draw[decorate, decoration={zigzag, segment length=5.5pt, amplitude=3.5pt}] (0,-2.56) -- (0,-3.52);
  \node[left=6pt] at (0,-3.04) {$k$};
  \draw (0,-3.52) -- (0,-3.80);

  % Bottom point and applied force F(t)
  \fill (0,-3.80) circle (2pt);
  \draw[->] (0,-3.80) -- (0,-4.48) node[pos=0.85, right=2pt] {$F(t)$};

  % Dimension ticks and arrows x1, x2
  \draw[thin] (0.68,-1.90) -- (1.45,-1.90);
  \draw[thin] (0.68,-3.80) -- (1.45,-3.80);
  \draw[<->] (1.15,0) -- (1.15,-1.90) node[midway, right=2pt] {$x_1$};
  \draw[<->] (1.15,-1.90) -- (1.15,-3.80) node[midway, right=2pt] {$x_2$};
\end{tikzpicture}

\end{center}
\end{problem}
\clearpage


\begin{problem}{2.1.61}
单位质量的质点在阻力与速度成正比的介质中运动, 比例系数为 $\beta$ , 在一个按指数衰减的力 $F_{0}e^{-\gamma t}$ (其中 $F_{0}, \gamma$ 均为常量) 作用下自静止开始运动. 用拉普拉斯变换法求解该质点此后的运动规律.
\end{problem}
\clearpage


\begin{problem}{2.1.62}
用积分因子法求解上题.
\end{problem}
\clearpage


\begin{problem}{2.1.63}
质量为 $m$ 的质点在图2.43所示的力作用下做直线运动.

(1) 求 $F(t)$ 的拉普拉斯变换的象函数 $f(p)$ ;

(2) 用拉普拉斯变换法求此运动的通解；

(3) 直接用逐次积分的方法求通解.

提示： $\frac{1}{p^{\mu}}\mathrm{e}^{-kp}(\mu >0)$ 的原函数为

\[
\frac {(t - k) ^ {\mu - 1}}{\Gamma (\mu)} u (t - k)
\]

其中

\[
\begin{array}{l} u (t - k) = \left\{ \begin{array}{l l} 0 & t <   k \\ 1 & t \geqslant k \end{array} \right. \\ \Gamma (\mu) = (\mu - 1)! \end{array}
\]
\begin{center}
\begin{tikzpicture}[>=Stealth, font=\small]
  \draw[->] (0,0) -- (4.8,0) node[below left=2pt and 0pt] {$t$};
  \draw[->] (0,0) -- (0,2.8) node[left=2pt] {$F$};
  \node[below left=1pt] at (0,0) {$O$};

  \draw[thick] (0,0) -- (1.4,2.0) -- (3.2,2.0);
  \draw[dashed] (0,2.0) node[left=2pt] {$F_0$} -- (1.4,2.0);
  \draw[dashed] (1.4,2.0) -- (1.4,0) node[below=2pt] {$t_1$};
  \draw[dashed] (3.2,2.0) -- (3.2,0) node[below=2pt] {$t_2$};
\end{tikzpicture}

\end{center}
\end{problem}
\clearpage


\begin{problem}{2.1.64}
质量为 m 的质点在恢复力 -kx 和阻力 $-\gamma \dot{x}$ 以及周期性驱动力(每个周期)

\[
F (t) = \left\{ \begin{array}{l l} A \sin \omega t & 0 <   t \leqslant \frac {\pi}{\omega} \\ 0 & \frac {\pi}{\omega} <   t \leqslant \frac {2 \pi}{\omega} \end{array} \right.
\]

的共同作用下做受迫振动. 用傅里叶级数求解稳态时的运动学方程(要求写出级数的前四项).
\end{problem}
\clearpage


\begin{problem}{2.1.65}
考虑一阻尼振子,其运动微分方程为

\[
\ddot {x} + \lambda \dot {x} | \dot {x} | + \omega_ {0} ^ {2} x = 0
\]

$\lambda$ 为小的正值常量, 初始条件为 $x(0) = a, \dot{x}(0) = 0$ . 用微扰法证明第一个后半周位移与前半周位移之比 (精确到 $\lambda$ 的一次方) 为 $1 - \frac{4}{3} \lambda a$ .
\end{problem}
\clearpage


\begin{problem}{2.1.66}
一物体在北纬 $40^{\circ}$ 的地球表面上方高 $h$ 处从静止降落, 若 $h = 100\mathrm{m}$ , 计算由于科里奥利力引起的着地点的横向位移.
\end{problem}
\clearpage


\begin{problem}{2.1.67}
由于地球的自转而引起的下列偏离的大小和方向是什么？

（1）从北纬 $\theta^{\circ}$ 、高 $L$ 的塔顶到底部悬挂的铅垂线下的摆锤；

(2) 从顶部降落的物体的着地点.
\end{problem}
\clearpage


\begin{problem}{2.1.68}
S 系是坐标原点固定在地球中心, z 轴指向北随地心做平动的惯性参考系, $S'$ 系与 S 系相似放置(原点与 z 轴完全重合)但是随地球自转的参考系.

(1) 写出任意矢量对时间的微商, 从 $S'$ 系到 $S$ 系的非相对论变换方程, 利用这变换方程推导出在 $S'$ 系中运动的物体受到的科里奥利力的表达式;

(2) 求出在北半球水平向东和垂直向上运动的物体受到的科里奥利力.
\end{problem}
\clearpage


\begin{problem}{2.1.69}
证明:一条宽度为 b、河水以速率 v 向北流的运河,在北纬 $\lambda$ 处,东岸的河水水面比西岸高 $\frac{2bv\omega\sin\lambda}{g}$ 。其中 a 是地球的自转角速度,g 为当地的表观重力加速度。
\end{problem}
\clearpage


\begin{problem}{2.1.70}
在北纬 $\lambda$ 处, 一个质量为 $m$ 的质点垂直上抛, 可达高度 $h$ 处. 问它在何处返回地面?
\end{problem}
\clearpage


\begin{problem}{2.1.71}
一个质点在地球北纬 $\lambda$ 处的一完全光滑的水平面(它与那个纬度的铅垂线垂直)上滑动.选 $z$ 轴铅直向上， $x$ 轴向南， $y$ 轴向东.若开始时质点位于 $x = 0, y = a, \dot{x} = u, \dot{y} = 0$ .

(1) 证明: 若略去运动微分方程中的 $\omega^2$ 项, 则坐在平面上的观察者看到的轨道方程为

\[
x ^ {2} + \left(y - a + \frac {u}{2 \omega \sin \lambda}\right) ^ {2} = \left(\frac {u}{2 \omega \sin \lambda}\right) ^ {2}
\]

(2) 说明略去 $\omega^{2}$ 项是不合理的.
\end{problem}
\clearpage


\begin{problem}{2.1.72}
当从正上方俯视时,发现在海洋表面之下有一个相当分明的孤立层中有一海洋环流做逆时针旋转,旋转周期是 14 小时.问此洋流是在什么纬度和哪个半球探测到的?
\end{problem}
\clearpage


\begin{problem}{2.1.73}
在北纬 $\lambda$ 处, 以初速 $v_{0}$ 向东发射炮弹, 发射角为 $\alpha$ . 若不计空气阻力, 视重力为常值, 且忽略其与铅垂线的偏差, 即认为重力指向地心. 求炮弹的运动方程以及着地点的偏差.
\end{problem}
\clearpage


\begin{problem}{2.1.74}
设有质量为 $m$ 的质点在椭球面

\[
\frac {x ^ {2}}{a ^ {2}} + \frac {y ^ {2}}{b ^ {2}} + \frac {(c - z) ^ {2}}{c ^ {2}} - 1 = 0
\]

的最低点做微振动， $z$ 轴竖直向上。求其运动规律及约束力。又如 $a = b$ ，且椭球绕 $z$ 轴以等角速 $\omega$ 转动，求质点在相对平衡状态时与椭球最低点的高度差。
\end{problem}
\clearpage


\begin{problem}{2.1.75}
一个质量为 $m$ 的质点，在重力作用下在一个光滑的旋转抛物面的内表面上运动。用柱坐标 $\rho, \varphi, z$ 表示，抛物面的方程为 $\rho^2 = az(a$ 为常量）， $z$ 轴竖直向上。开始质点于 $\rho = a, \varphi = 0, z = a$ 处被抛射时，具有在水平方向的速率 $v$ 。求质点运动过程中 $z$ 的最大值和最小值。
\end{problem}
\clearpage


\begin{problem}{2.1.76}
一个单位质量的质点在一个旋转抛物面内表面上运动,旋转抛物面的柱坐标方程为 $\rho^2 = 4az, z$ 轴竖直向上， $a$ 为常量。该质点受到一个垂直于 $z$ 轴向外的大小为 $\mu \rho$ 的斥力，摩擦力可忽略不计，开始质点位于 $z = 0$ ，具有垂直于 $z$ 轴的速度 $2a\mu^{1/2}$ 。用拉格朗日乘子法求质点所能达到的最高点。
\end{problem}
\clearpage


\begin{problem}{2.1.77}
一块均匀的面密度为 $\sigma$ 的边长为 $2a$ 和 $2b$ 的矩形平板, 一个质量为 $m$ 的质点位于通过矩形板的中心且垂直于平板的直线上的一点, 质点距平板的距离为 $z$ . 求矩形

平板作用于质点的万有引力.
\end{problem}
\clearpage


\begin{problem}{2.1.78}
求一个质量为 $M$ 的均质球作用于位于离球心距离为 $r$ 的质量为 $m$ 的质点的万有引力. 设球的半径为 $R$ .
\end{problem}
\clearpage


\begin{problem}{2.1.82}
考虑一块靠自身引力维持的由液态物质构成并处于流体静力平衡的平板，其总厚度为 $2h$ ，横向范围 $(x$ 和 $y$ 方向)是无限的. 平板密度只是 $z$ 的函数 $\rho(z)$ ，与 $x, y$ 无关，且对于 $z = 0$ 的中间平面是对称的. 求中间平面的压强 $p(0)$ ，把它写成 $\sigma = \int_{0}^{h} \rho(z) \mathrm{d}z$ 的函数.
\end{problem}
\clearpage


\section{质点的动能定理和机械能守恒定律}

\begin{problem}{2.2.1}
一质点在力 $F = 4yi + 2xj + k$ 作用下沿一螺旋线 $x = 4\cos \theta, y = 4\sin \theta, z = 2\theta$ 从 $\theta = 0$ 到 $\theta = 2\pi$ . 求此力在这过程中对质点所做的功.
\end{problem}
\clearpage


\begin{problem}{2.2.2}
作用在骑自行车的人身上的摩擦力和空气阻力合在一起,其大小 F=av,其中 v 是骑车人的速度, 而 $a=4N \cdot s/m$ , 该骑车人最多能产生 600W 的推动功率. 问在无风情况下在水平面上他的最大速度等于多少?
\end{problem}
\clearpage


\begin{problem}{2.2.3}
一辆 2500kg 的卡车在水平面上以 30m/s 的速度行驶. 现突然推入空挡(即做滑行), 于是速度按下列规律变化:

\[
v = \frac {3 0}{1 + \frac {t}{3 0}} \quad (\mathrm{m/s})
\]

其中 t 以秒作为单位. 求在这条路上以 15m/s 的速度驱动该卡车所需要的功率.
\end{problem}
\clearpage


\begin{problem}{2.2.4}
一个无限大的均质薄板割出一个半径为 $a$ 的圆孔，一个半径为 $a$ 的均质球壳放在孔中，圆孔中心和球壳中心重合。薄板和球壳单位面积的质量均为 $\sigma$ 。要将质量为 $m$ 的粒子从孔的中心 $O$ 点移到在过 $O$ 垂直于薄板的直线上离 $O$ 点距离为 $h(h \geqslant a)$ 的 $P$ 点，至少要做多大的功。
\end{problem}
\clearpage


\begin{problem}{2.2.5}
一重量为 w 的人处于一个重量为 w 的电梯内, 电梯以加速度 a 向上加速, 在某一瞬时的速度为 V.

(1) 问人的表观重量等于多少?

(2) 若此人以相对于电梯为 v 的速度爬上安放在电梯内的一个垂直梯子, 此人的能量消耗率(功率输出)等于多少?
\end{problem}
\clearpage


\begin{problem}{2.2.6}
一质量为 m 的质点, 能在半径为 a 的圆环形小管内无摩擦滑动, 圆环绕铅直直径以恒定角速度 $\omega$ 转动, 如图 2.51 所示. 写出质点的运动微分方程; 如果质点在 $\theta=0$ 的不稳定平衡位置, 受轻微扰动, 求出其最大动能位置.
\begin{center}
\begin{tikzpicture}[>=Stealth, font=\small]
  \def\R{1.85}
  \def\r{1.55}
  \def\rmid{1.70}
  \def\ang{45}

  % Concentric circular tube
  \draw[thick] (0,0) circle (\R);
  \draw[thick] (0,0) circle (\r);

  % Vertical rotation axis
  \draw[thick] (0,-2.35) -- (0,2.65);

  % Rotation arrow and omega
  \draw[->, thick] (-0.25,2.42) arc (145:405:0.42 and 0.14);
  \node[right=2pt] at (0.38,2.55) {$\omega$};

  % Gravity vector on the left
  \draw[->, thick] (-2.45,0.65) -- (-2.45,-0.65);
  \node[left=2pt] at (-2.45,0) {$g$};

  % Radial line to mass m
  \coordinate (O) at (0,0);
  \coordinate (Top) at (0,1);
  \coordinate (M) at ({90-\ang}:\r);
  \draw[thick] (O) -- (M);
  \fill (M) circle (2.2pt);
  \node at ({90-\ang}:{\R+0.32}) {$m$};

  % Angle theta
  \pic[draw, angle radius=0.65cm, angle eccentricity=1.35, "$\theta$"] {angle=M--O--Top};
\end{tikzpicture}

\end{center}
\end{problem}
\clearpage


\begin{problem}{2.2.7}
找出与下列势能函数相应的力场：

(1) $V = \frac{1}{2}\ln (x^2 +y^2 +z^2);$ 

(2) $V(r, \theta, \varphi) = \frac{1}{r} \mathrm{e}^{-kr}$ ; 

(3) $V(r, \theta, \varphi) = \frac{a}{r^2} \cos \theta$ .
\end{problem}
\clearpage


\begin{problem}{2.2.8}
证明一个半径为 $a$ 、面密度为 $\sigma$ 的薄圆盘，圆盘的轴向距离为 $x$ 的点的引力势为

\[
V (x) = - 2 \pi G \sigma (\sqrt {x ^ {2} + a ^ {2}} - | x |)
\]
\end{problem}
\clearpage


\begin{problem}{2.2.9}
计算半径为 a、质量为 M 的均质圆环对位于环的平面内，但在环的外侧的各点的引力势.

提示: 可表示为第一类椭圆积分.
\begin{center}
\begin{tikzpicture}[>=Stealth, font=\small]
  \def\a{1.55}
  \def\x{4.1}
  \def\ang{42}

  \coordinate (O) at (0,0);
  \coordinate (A) at (\ang:\a);
  \coordinate (P) at (\x,0);

  % Ring
  \draw[thick] (O) circle (\a);

  % Center O
  \fill (O) circle (1.8pt);
  \node[below=2pt] at (O) {$O$};

  % Triangle sides
  \draw[thick] (O) -- (A) node[midway, above left=0pt] {$a$};
  \draw[thick] (O) -- (P) node[pos=0.68, below=2pt] {$x$};
  \draw[thick] (A) -- (P);

  % Point P
  \node[right=2pt] at (P) {$P$};

  % Angle theta
  \pic[draw, angle radius=0.58cm, angle eccentricity=1.38, "$\theta$"] {angle=P--O--A};
\end{tikzpicture}

\end{center}
\end{problem}
\clearpage


\begin{problem}{2.2.10}
考虑一个任意形状的物体,证明:对物体外一个球内各点的引力势的平均值等于在此球心处的引力势.
\end{problem}
\clearpage


\begin{problem}{2.2.11}
判断以下各力是否是保守力,如果是保守力,找出它的势能函数.

(1) $F=(ax+by^{2})i+(az+2bxy)j+(ay+bz^{2})k,a,b$ 为常量；

(2) $F=a\times r,a$ 为常矢量；

(3) $F = ar, a$ 为常量；

(4) $F=a(a\cdot r)$ ,a为常矢量.
\end{problem}
\clearpage


\begin{problem}{2.2.12}
一个密度均匀的行星绕一个固定轴以角速度 $\omega$ 自转，行星的赤道半径 $R_{\mathrm{E}}$ 略大于球的半径 $R$ ，行星两极的半径 $R_{\mathrm{P}}$ 略小于球的半径 $R$ ，用参量 $\varepsilon = \frac{R_{\mathrm{E}} - R_{\mathrm{P}}}{R_{\mathrm{E}}}$ 描述变形的大小.由于这种变形引起重力势有一增量

\[
\Delta V (R, \theta) = - \frac {2}{5} \frac {G M \varepsilon R _ {\mathrm{E}} ^ {2}}{R ^ {3}} P _ {2} (\cos \theta)
\]

其中 $\theta$ 是球坐标， $P_{2}(\cos \theta) = \frac{3}{2}\cos^{2}\theta - \frac{1}{2}$ . 由行星表面平衡的条件得出 $\varepsilon$ 与 $\lambda = \frac{\omega^2 R_E}{g}$ 的关系，对地球就 $\varepsilon$ 作数值估计.
\end{problem}
\clearpage


\begin{problem}{2.2.13}
求在势能为

\[
V (x) = \frac {1}{2} k x ^ {2} - \frac {1}{3} m \lambda x ^ {3}
\]

(其中 $\lambda$ 很小) 的保守力作用下的非线性振子的运动的一级近似解, 假定 t=0 时 x=0.
\end{problem}
\clearpage


\begin{problem}{2.2.14}
(1) 质量为 $m$ 的粒子在势能 $V(x) = \frac{cx}{x^2 + a^2}$ 的保守力场中运动, 其中 $c$ 和 $a$ 是正的常数. 求稳定平衡位置及在它附近做小振动的周期.

(2) 如果粒子从所求的平衡位置以速度 v 开始运动, 求 v 值的范围, 粒子将做: (a) 振动; (b) 逃到 $-\infty$ ; (c) 逃到 $+\infty$ .
\end{problem}
\clearpage


\begin{problem}{2.2.15}
比地面高 h 的平台上有一质量为 m 的小车系于绳子一端, 绳子跨过一小滑轮在另一端被地面上的人拉着并以匀速 $v_{0}$ 向右运动，如图2.58所示.若小车与平台间的摩擦力可以不计，求：

(1) 当人位于离小滑轮正下方向右距离 s 处小车的速度和加速度；

(2) 人在上述位置拉绳的力；

(3) 在人从滑轮正下方到上述位置的过程中人所做的功.
\begin{center}
\begin{tikzpicture}[>=Stealth, font=\small]
  % Platform and lower ground
  \draw[thick] (-3.1,1.65) -- (-0.25,1.65) arc (90:0:0.25) -- (0,0) -- (3.8,0);

  % Hatching for platform, wall, and ground
  \foreach \x in {-2.95,-2.65,...,-0.25} {
    \draw (\x,1.65) -- (\x-0.2,1.42);
  }
  \foreach \y in {0.2,0.5,0.8,1.1,1.4} {
    \draw (0,\y) -- (-0.2,\y-0.2);
  }
  \foreach \x in {0.2,0.5,...,3.6} {
    \draw (\x,0) -- (\x-0.2,-0.23);
  }

  % Pulley
  \coordinate (Pul) at (0.22,1.92);
  \def\rp{0.24}
  \draw[thick] (Pul) circle (\rp);
  \fill (Pul) circle (1.5pt);

  % Cart
  \draw[thick] (-2.6,1.93) rectangle (-1.5,2.52);
  \draw[thick] (-2.35,1.79) circle (0.14);
  \draw[thick] (-1.75,1.79) circle (0.14);

  % Rope from cart to pulley top
  \pgfmathsetmacro{\yrope}{1.92+\rp}
  \draw[thick] (-1.5,\yrope) -- (0.22,\yrope);

  % Exact tangent from person's hand (2.68,0.52) to pulley
  \coordinate (Hand) at (2.68,0.52);
  \pgfmathsetmacro{\dx}{2.68-0.22}
  \pgfmathsetmacro{\dy}{0.52-1.92}
  \pgfmathsetmacro{\dist}{sqrt(\dx*\dx+\dy*\dy)}
  \pgfmathsetmacro{\baseang}{atan2(\dy,\dx)}
  \pgfmathsetmacro{\offsetang}{acos(\rp/\dist)}
  \pgfmathsetmacro{\tang}{\baseang+\offsetang}
  \coordinate (TanPt) at ($(Pul)+(\tang:\rp)$);
  \draw[thick] (TanPt) -- (Hand);

  % Height dimension h
  \draw (-1.0,0) -- (-0.2,0);
  \draw[<->] (-0.8,0) -- (-0.8,\yrope) node[midway, left=2pt] {$h$};

  % Person pulling the rope
  \draw[thick] (2.95,0.78) circle (0.11); % head
  \draw[thick] (2.95,0.67) -- (2.92,0.34); % torso
  \draw[thick] (2.94,0.58) -- (Hand); % back arm holding rope
  \draw[thick] (2.94,0.58) -- (3.08,0.46) -- (3.12,0.57); % front arm
  \draw[thick] (2.92,0.34) -- (2.78,0.18) -- (2.72,0); % back leg
  \draw[thick] (2.92,0.34) -- (3.06,0.18) -- (2.98,0); % front leg

  % Velocity arrow v_0
  \draw[->, thick] (3.25,0.50) -- (3.80,0.50) node[above=2pt, pos=0.65] {$v_0$};
\end{tikzpicture}

\end{center}
\end{problem}
\clearpage


\begin{problem}{2.2.16}
一质量为 $m$ 的小球系于不可伸长的轻绳的一端，穿过桌面上的一小孔后在光滑的水平桌面上运动，绳子与小孔间无摩擦力。问：

(1) 要使小球以小孔为中心做半径为 $r_{1}$ 、角速度为 $\omega_{1}$ 的圆周运动，绳的另一端需用多大的拉力？

(2) 要使小球圆周运动的半径从 $r=r_{1}$ 缩小到 $r=r_{2}$ ，拉力需做多大的功？
\end{problem}
\clearpage


\begin{problem}{2.2.17}
火车质量为 $m$ ，其所产生功率为常数 $P$ ，若车受的阻力为常数 $f$ ，试证时间与速度的关系为 $t = \frac{mP}{f^2} \ln \frac{P}{P - fv} - \frac{mv}{f}$ ；若所受阻力与速度 $v$ 成正比，则 $t = \frac{mv}{2f} \ln \frac{P}{P - fv}$ .
\end{problem}
\clearpage


\begin{problem}{2.2.18}
一质量为 $m$ 的质点, 受引力作用在一直线上运动. 当 $x \geqslant a$ 时, 引力值为 $m\mu a^2 / x^2$ ; 当 $x \leqslant a$ 时, 引力值为 $m\mu x / a$ , 式中 $x$ 是相对于线上某一点 (取为原点) 的距离, $\mu$ 及 $a$ 为常量. 如质点在离原点 $2a$ 处从静止开始运动, 求到达原点时的速度及走此段路程所需的时间.
\end{problem}
\clearpage


\begin{problem}{2.2.19}
轴套 A 的质量为 0.5kg，沿着位于铅垂平面的螺线形杆无摩擦地滑动。螺线形方程为 $r=0.3\theta$ ，式中 r 以米计， $\theta$ 以弧度计，轴套在大小不变的径向力 T=10N 作用下，从 A 点静止释放，滑到位置 B 处（见图 2.60）。求轴套到达 B 处时的速度。
\begin{center}
\begin{tikzpicture}[>=Stealth, font=\small]
  \def\c{1.1}
  \def\w{0.055}

  % Axes
  \draw (-1.3,0) -- (4.1,0);
  \draw (0,-0.55) -- (0,2.75);
  \node[below left=1pt] at (0,0) {$O$};

  % Spiral rod center path: x(t) = -c*t*cos(t), y(t) = c*t*sin(t)
  % Draw inner and outer boundaries as a smooth double-walled tube
  \draw[thick, double=white, double distance=2.6pt, samples=120, domain=0:190]
    plot ({-\c*(\x*pi/180)*cos(\x)}, {\c*(\x*pi/180)*sin(\x)});

  % Hatching support at O
  \draw[thick] (0,-0.35) -- (0,0.35);
  \foreach \y in {-0.25,-0.05,0.15,0.32} {
    \draw (0,\y) -- (0.16,\y-0.16);
  }

  % Hatching base below B
  \draw[thick] (3.3,-0.31) -- (4.05,-0.31);
  \foreach \x in {3.42,3.62,3.82,4.02} {
    \draw (\x,-0.31) -- (\x-0.14,-0.47);
  }

  % Point A (theta = 90 deg) and Point B (theta = 180 deg)
  \pgfmathsetmacro{\yA}{\c*pi/2}
  \pgfmathsetmacro{\xB}{\c*pi}
  \node[above left=2pt] at (0,\yA) {$A$};
  \node[above right=2pt] at (\xB+0.05,0.02) {$B$};

  % Radial segment r at theta = 62 deg
  \def\thOne{62}
  \pgfmathsetmacro{\rOne}{\c*\thOne*pi/180}
  \coordinate (O) at (0,0);
  \coordinate (NegX) at (-1,0);
  \coordinate (Pone) at ({-\rOne*cos(\thOne)}, {\rOne*sin(\thOne)});
  \draw[thick] (O) -- (Pone) node[pos=0.52, right=1pt] {$r$};
  \pic[draw, angle radius=0.68cm, angle eccentricity=1.28, "$\theta$"] {angle=Pone--O--NegX};

  % Collar / sleeve and radial force T at theta = 142 deg
  \def\thTwo{142}
  \pgfmathsetmacro{\rTwo}{\c*\thTwo*pi/180}
  \pgfmathsetmacro{\px}{-\rTwo*cos(\thTwo)}
  \pgfmathsetmacro{\py}{\rTwo*sin(\thTwo)}
  % Tangent angle of spiral at theta = 142 deg
  \pgfmathsetmacro{\dx}{-\c*(cos(\thTwo)-(\thTwo*pi/180)*sin(\thTwo))}
  \pgfmathsetmacro{\dy}{\c*(sin(\thTwo)+(\thTwo*pi/180)*cos(\thTwo))}
  \pgfmathsetmacro{\tang}{atan2(\dy,\dx)}

  \coordinate (Ptwo) at (\px,\py);
  \coordinate (Ttip) at ({\px+0.85*(-cos(\thTwo))}, {\py+0.85*sin(\thTwo)});
  \draw[dashed, thick] (O) -- (Ptwo);

  % Sleeve around the spiral rod at Ptwo
  \begin{scope}[shift={(Ptwo)}, rotate=\tang]
    \draw[thick, fill=white] (-0.11,-0.09) rectangle (0.11,0.09);
  \end{scope}

  \draw[->, thick] (Ptwo) -- (Ttip) node[above right=0pt] {$\boldsymbol{T}$};
\end{tikzpicture}

\end{center}
\end{problem}
\clearpage


\begin{problem}{2.2.20}
一质量为 80kg 的人由 1m 高处跳下, 当他落地时, 他忘了弯曲他的膝盖, 人体减速的距离只有 1cm. 求在减速阶段, 作用在他腿上的力.
\end{problem}
\clearpage


\begin{problem}{2.2.21}
一质点无摩擦地在环形轨道上滑下,如图 2.61 所示,轨道弯曲段的曲率半径为 R,该质点由高 h 处由自静止开始下滑,在某处质点开始和轨道脱离接触,说明脱离接触的位置并计算出发生这种情况的 h 的最小值.
\begin{center}
\begin{tikzpicture}[>=Stealth, font=\small]
  \def\R{1.45}
  \pgfmathsetmacro{\xCtwo}{\R*sqrt(3)}
  \pgfmathsetmacro{\xA}{\R*sqrt(3)/2}
  \pgfmathsetmacro{\yA}{\R/2}

  \coordinate (C1) at (0,\R);
  \coordinate (C2) at (\xCtwo,0);
  \coordinate (A) at (\xA,\yA);
  \coordinate (B) at (\xCtwo,\R);

  % Ground
  \draw[thick] (-2.85,0) -- (3.85,0);
  \foreach \x in {-2.65,-2.35,...,3.75} {
    \draw (\x,0) -- (\x-0.22,-0.25);
  }

  % Left incline + valley arc (225 deg to 330 deg around C1) + hill arc (150 deg down to 35 deg around C2)
  \coordinate (JoinL) at ($(C1)+(225:\R)$);
  \coordinate (TopL) at ($(JoinL)+(-1.55,1.55)$);
  \draw[thick] (TopL) -- (JoinL)
    arc (225:330:\R)
    arc (150:35:\R);

  % Block near top of left incline (rotated -45 deg)
  \coordinate (Blk) at ($(TopL)+(0.16,-0.16)$);
  \begin{scope}[shift={(Blk)}, rotate=-45]
    \draw[thick, fill=white] (-0.18,0) rectangle (0.18,0.22);
    \draw (-0.06,0) -- (-0.06,0.22) (0.06,0) -- (0.06,0.22);
  \end{scope}

  % Height h arrow
  \draw[<->] (Blk |- 0,0) -- (Blk) node[midway, left=2pt] {$h$};

  % Centers, dashed construction lines, and radius arrows
  \fill (C1) circle (1.5pt);
  \fill (C2) circle (1.5pt);
  \draw[dashed] (C1) -- (B);
  \draw[dashed] (C2) -- (A);

  % Radius arrows labeled R
  \draw[->, thick] (C1) -- ($(C1)+(302:\R)$) node[midway, left=1pt] {$R$};
  \draw[->, thick] (C2) -- ($(C2)+(110:\R)$) node[midway, right=1pt] {$R$};

  % Angle 30 deg at C2
  \coordinate (LeftG) at (0,0);
  \pic[draw, angle radius=0.65cm, angle eccentricity=1.55, "$30^\circ$"] {angle=A--C2--LeftG};

  % Labels A and B
  \node[above left=1pt] at (A) {$A$};
  \node[above=2pt] at (B) {$B$};
\end{tikzpicture}

\end{center}
\end{problem}
\clearpage


\begin{problem}{2.2.22}
一个转动的球形行星, 其赤道上的点的速率为 $V$ , 赤道上的 $g$ 是两极处 $g$ 的一半. 问粒子从极点逃逸的逃逸速度多大?
\end{problem}
\clearpage


\begin{problem}{2.2.23}
一跳伞员在 3000m 高处起跳. 在降落伞张开之前他达到一极限速度 30m/s.

(1) 假定空气阻力正比于速度,他达到这个极限速度需多久?达到这个速度时,运行了多少距离?

(2) 为了经受一个不大于 $10g(g$ 为重力加速度)的减速度, 在他的伞张开以后, 当他碰到地面时, 速度已降到 $3m/s$ , 他弯曲膝部以减缓冲撞, 需将膝部弯曲多少? 假定他的膝部像弹簧一样具有正比于位移的阻力.

(3) 空气阻力正比于速度的假定是否合理?
\end{problem}
\clearpage


\begin{problem}{2.2.24}
一个光滑球固定在水平面上, 一个点粒子在最高点从静止开始沿球面滑下. 令球的半径为 R, 描述粒子落到平面以前的路径 (见图 2.62).
\begin{center}
\begin{tikzpicture}[>=Stealth, font=\small]
  \def\R{1.5}
  \def\th{40}

  % Ground / x-axis
  \draw[->, thick] (-2.0,0) -- (4.0,0) node[below=2pt] {$x$};
  \foreach \x in {-1.8,-1.5,...,3.6} {
    \draw (\x,0) -- (\x-0.22,-0.24);
  }

  % Sphere
  \coordinate (C) at (0,\R);
  \draw[thick] (C) circle (\R);
  \fill (C) circle (1.6pt);

  % y-axis
  \draw[->, thick] (0,0) -- (0,3.6) node[left=2pt] {$y$};
  \node[left=2pt] at (0,2.35) {$R$};

  % Particle on sphere at angle th from top vertical
  \coordinate (Top) at (0,2*\R);
  \coordinate (P) at ($(C)+({90-\th}:\R)$);
  \draw[thick] (C) -- (P);
  \fill (P) circle (2pt);

  % Angle theta
  \pic[draw, angle radius=0.55cm, angle eccentricity=1.35, "$\theta$"] {angle=P--C--Top};

  % Force & velocity arrows at P
  \draw[->, thick] (P) -- ($(P)+({90-\th}:0.62)$) node[right=1pt] {$N$};
  \draw[->, thick] (P) -- ($(P)+(0,-0.72)$) node[below=1pt] {$mg$};
  \draw[->, thick] (P) -- ($(P)+(-\th:0.75)$) node[above right=0pt] {$v$};

  % Parabolic trajectory after leaving sphere
  \draw[dashed, thick] (P) .. controls ($(P)+(-\th:0.95)$) and ($(P)+(1.45,-1.25)$) .. (3.15,0);
\end{tikzpicture}

\end{center}
\end{problem}
\clearpage


\begin{problem}{2.2.25}
若在北京和天津间用一条直的地下铁道联结, 如图 2.63 所示, 两市间的火车在地球引力作用下运行, 两市间直线距离为 $130\mathrm{km}$ , 地球半径 $R = 6400\mathrm{km}$ . 忽略摩擦, 计算火车的最大速度以及从北京到天津坐此火车所需的时间.
\begin{center}
\begin{tikzpicture}[>=Stealth, font=\small]
  \def\R{3.2}
  \def\aEnd{52}
  \def\aChord{58}

  \coordinate (O) at (0,0);
  \fill (O) circle (1.6pt);

  % Earth surface arc
  \draw[very thick] (\aEnd:\R) arc (\aEnd:{180-\aEnd}:\R);

  % Straight tunnel chord
  \coordinate (P1) at ({180-\aChord}:\R);
  \coordinate (P2) at (\aChord:\R);
  \draw[thick] (P1) -- (P2);

  % Dimension 130km
  \pgfmathsetmacro{\yCh}{\R*sin(\aChord)}
  \pgfmathsetmacro{\xCh}{\R*cos(\aChord)}
  \draw (-\xCh,\yCh) -- (-\xCh,\yCh-0.56);
  \draw (\xCh,\yCh) -- (\xCh,\yCh-0.56);
  \draw[<->] (-\xCh,\yCh-0.42) -- (\xCh,\yCh-0.42) node[midway, above=1pt] {$130\,\mathrm{km}$};

  % Radius arrow R
  \draw[->, thick] (O) -- (\aEnd:\R) node[midway, below right=1pt] {$R$};
\end{tikzpicture}

\end{center}
\end{problem}
\clearpage


\begin{problem}{2.2.26}
估计一下通过跳跃就能脱离的小行星多大.
\end{problem}
\clearpage


\begin{problem}{2.2.27}
求:(1)月球表面的重力加速度;(2)从月球表面出发逃离月球所需的最小速度.月球质量为 $7.35\times10^{22}kg$ ,月球半径为 $1.74\times10^{3}km$ .
\end{problem}
\clearpage


\begin{problem}{2.2.28}
一质量为 m 的质点从置于光滑水平面上的、质量为 M 的光滑的、半径为 R 的弧形槽的顶端滑下，如图 2.65 所示。开始滑下时，m、M 都是静止的，求质点离开弧形槽时质点和弧形槽的速度。
\begin{center}
\begin{tikzpicture}[>=Stealth, font=\small]
  \def\R{2.1}
  \def\hbase{0.35}
  \def\topY{2.45} % hbase + R

  % Ground
  \draw[thick] (-0.4,0) -- (4.1,0);
  \foreach \x in {-0.2,0.1,...,3.9} {
    \draw (\x,0) -- (\x-0.2,-0.22);
  }

  % Block M with quarter-circle track
  % Center of arc at (3.1, 2.45), left top lip at (1.0, 2.45), bottom lip at (3.1, 0.35)
  \draw[thick] (0.35,0) -- (0.35,\topY) -- (1.0,\topY)
    arc (180:270:\R) -- (3.1,0) -- cycle;
  % Base step line
  \draw (0.35,0.1) -- (3.1,0.1);

  % Center of curvature and dashed lines
  \coordinate (C) at (3.1,\topY);
  \fill (C) circle (1.5pt);
  \draw[dashed, thick] (1.0,\topY) -- (C) -- (3.1,\hbase);

  % Right-angle mark at top lip
  \draw (1.0,{\topY-0.16}) -- (1.16,{\topY-0.16}) -- (1.16,\topY);

  % Radius arrow R
  \draw[->, thick] (C) -- ($(C)+(225:\R)$) node[midway, above left=0pt] {$R$};

  % Labels m and M
  \node[above=2pt] at (1.55,\topY) {$m$};
  \node at (0.98,0.85) {$M$};
\end{tikzpicture}

\end{center}
\end{problem}
\clearpage


\begin{problem}{2.2.29}
如图 2.66 所示, 小球质量为 m, 连接在绕水平轴 O 旋转的 BD 杆的 D 端, 连杆 BC 穿过绕 A 点转动的套筒, 并能在其内滑动, 劲度系数为 k 的弹簧套在 BC 杆的 AC 段, 当球下降时压缩弹簧, $\theta = 60^{\circ}$ 的位置时, 弹簧无伸缩. 从此位置静止释放, 求 $\theta = 90^{\circ}$ 时 BC 杆的角速度. 设 BC 杆及 BD 杆的质量及各处摩擦可略去不计.
\begin{center}
\begin{tikzpicture}[>=Stealth, font=\small]
  \def\b{2.1}
  \def\th{65}

  \coordinate (O) at (0,0);
  \coordinate (A) at (0,-\b);
  \coordinate (B) at ({-\b*sin(\th)}, {-\b*cos(\th)});
  \coordinate (D) at ({\b*sin(\th)}, {\b*cos(\th)});

  % Ceiling support at O
  \draw[thick] (-0.36,0.45) -- (0.36,0.45);
  \foreach \x in {-0.30,-0.15,0.0,0.15,0.30} {
    \draw (\x,0.45) -- (\x+0.12,0.59);
  }
  \draw[thick] (-0.18,0.45) -- (-0.07,0.07);
  \draw[thick] (0.18,0.45) -- (0.07,0.07);

  % Rod BD
  \draw[thick, double=white, double distance=2.6pt] (B) -- (D);
  \node[above left=1pt] at ($(O)!0.52!(B)$) {$b$};
  \node[above left=1pt] at ($(O)!0.52!(D)$) {$b$};

  % Mass sphere at D
  \draw[thick, fill=white] (D) circle (0.24);
  \node[above left=3pt and 4pt] at (D) {$D$};

  % Vertical dashed line OA and dimension b
  \draw[dashed, thick] (O) -- (A);
  \draw (0.22,0) -- (0.72,0);
  \draw (0.32,-\b) -- (0.72,-\b);
  \draw[<->] (0.56,0) -- (0.56,-\b) node[midway, right=2pt] {$b$};

  % Angle theta
  \pic[draw, angle radius=0.45cm, angle eccentricity=1.45, "$\theta$"] {angle=B--O--A};

  % Rod BC direction from B through A to C
  \pgfmathsetmacro{\dx}{0 - (-\b*sin(\th))}
  \pgfmathsetmacro{\dy}{-\b - (-\b*cos(\th))}
  \pgfmathsetmacro{\angBA}{atan2(\dy,\dx)}
  \coordinate (C) at ($(A)+(\angBA:2.25)$);

  % Rod BC (double line)
  \draw[thick, double=white, double distance=2.6pt] (B) -- (C);

  % Spring along AC (from sleeve exit to plate C)
  \coordinate (SpStart) at ($(A)+(\angBA:0.32)$);
  \coordinate (SpEnd) at ($(C)+(\angBA:-0.08)$);
  \draw[thick, decorate, decoration={coil, aspect=0.45, segment length=2.8mm, amplitude=2.1mm}]
    (SpStart) -- (SpEnd);
  \node[above right=5pt and 2pt] at ($(SpStart)!0.5!(SpEnd)$) {$k$};

  % End plate at C
  \begin{scope}[shift={(C)}, rotate=\angBA]
    \draw[thick, fill=white] (-0.06,-0.22) rectangle (0.06,0.22);
  \end{scope}
  \node[right=4pt] at (C) {$C$};

  % Sleeve at A
  \begin{scope}[shift={(A)}, rotate=\angBA]
    \draw[thick, fill=white] (-0.32,-0.15) rectangle (0.32,0.15);
  \end{scope}

  % Floor support at A
  \draw[thick] (-0.42,-\b-0.48) -- (0.42,-\b-0.48);
  \foreach \x in {-0.35,-0.20,-0.05,0.10,0.25,0.38} {
    \draw (\x,-\b-0.48) -- (\x-0.12,-\b-0.62);
  }
  \draw[thick] (-0.24,-\b-0.48) -- (-0.08,-\b-0.14);
  \draw[thick] (0.24,-\b-0.48) -- (0.08,-\b-0.14);

  % Pivot pins at O, B, A
  \draw[thick, fill=white] (O) circle (2.4pt);
  \draw[thick, fill=white] (B) circle (2.4pt);
  \draw[thick, fill=white] (A) circle (2.0pt);

  \node[above left=2pt and 6pt] at (O) {$O$};
  \node[left=3pt] at (B) {$B$};
  \node[below left=2pt and 7pt] at (A) {$A$};
\end{tikzpicture}

\end{center}
\end{problem}
\clearpage


\begin{problem}{2.2.30}
求质量为 $m$ 和具有能量 $E$ 的质点在势场 $V = V_{0}\tan^{2}ax$ (式中 $V_{0}$ 和 $a$ 为常量) 中做一维运动的周期.
\end{problem}
\clearpage


\begin{problem}{2.2.31}
一根长度为 $a$ 的不可伸长的轻绳，一端连着一个质量为 $m$ 的粒子 $P$ ，另一端连着一个质量为 $\frac{3}{2} m$ 的环 $Q$ ，此环穿在一根固定的粗糙的水平的金属丝上。从环和粒子处于平衡的状态给粒子一个沿水平方向的速度 $\sqrt{2ag}$ ，证明：在小环不在金属丝上滑动的情况下，绳中张力为 $3mg\cos \theta$ ，其中 $\theta$ 是绳子与铅垂线间的夹角；并证明：如环与金属丝间摩擦因数大于 $\frac{1}{\sqrt{3}}$ ，环就不会滑动。
\end{problem}
\clearpage


\begin{problem}{2.2.32}
一质点自一光滑的抛物线上一点 $Q$ 自由滑下, 该点在轴线上方 $h$ 处, 抛物线的方程为 $y^{2} = 2px$ , $y$ 轴竖直向上. 问质点滑至何处时, 曲线对质点的约束力将改变符号 (求出此点的 $y$ 坐标满足的方程即可).
\end{problem}
\clearpage


\begin{problem}{2.2.33}
一个质点在恒定的重力作用下在 $xy$ 平面内运动, 重力沿 $y$ 轴负方向. 找出该二维运动的四个独立的运动常数, 其中有三个独立的运动常数不显含时间.
\end{problem}
\clearpage


\begin{problem}{2.2.34}
当一个正在收缩的星体收缩到它的 Schwarzschild 半径 $R_{s}$ 以下时, 它变成一个黑洞, 光和任何粒子都不能逃离它. 导出 $R_{s}$ 和它的质量 M 的关系.
\end{problem}
\clearpage


\section{质点的角动量定理和角动量守恒定律}

\begin{problem}{2.3.1}
一质量为 m 的质点在光滑水平面上以速率 $v_{0}$ 做半径为 $R_{0}$ 的圆周运动。该质点系在一根不可伸长的轻绳上，绳子又穿过该平面上的一个光滑小孔，如图 2.69 所示。

(1) 绳中的张力多大?

(2) 质点对小孔的角动量多大?

(3) 质点的动能多大?

(4) 若使绳中的张力逐渐地增大, 最后使质点做

半径为 $\frac{1}{2} R_0$ 的圆周运动，质点最终的动能多大？

(5) 如果张力不是逐渐增大, 上问的答案还对吗?
\begin{center}
\begin{tikzpicture}[>=Stealth, thick]
  % Horizontal plane (top face and front/right thickness)
  \draw[fill=white] (-2.2,-1.1) -- (1.6,-1.1) -- (2.8,1.1) -- (-1.0,1.1) -- cycle;
  \draw (-2.2,-1.1) -- (-2.2,-1.24) -- (1.6,-1.24) -- (1.6,-1.1);
  \draw (1.6,-1.24) -- (2.8,0.96) -- (2.8,1.1);

  % Circular orbit in perspective
  \coordinate (O) at (0.3,0.02);
  \draw (O) ellipse [x radius=1.35, y radius=0.52];
  \fill (O) circle (1.4pt);

  % Point on orbit, radius string, and tangential velocity
  \coordinate (P) at ($(O) + (35:1.35 and 0.52)$);
  \draw (O) -- (P);
  \draw[->] (P) -- ++(-0.78, 0.45) node[above right=-1pt] {$v_0$};

  % Vertical string passing through the hole
  \draw[dashed] (O) -- (0.3,-1.24);
  \draw (0.3,-1.24) -- (0.3,-2.25);
\end{tikzpicture}

\end{center}
\end{problem}
\clearpage


\begin{problem}{2.3.2}
一质量为 $m$ 的小孩坐在一秋千上, 秋千的质量可以忽略, 悬挂在一端固定的、长度为 $l$ 的绳的另一端. 小孩的父亲拉他的后背使绳子和铅直方向的夹角为 $1\mathrm{rad}$ , 然后以力 $F = mg$ 沿圆周的切线方向推小孩直到绳子到达铅直位置放开该秋千. 问孩子的父亲推秋千用了多少时间? 在 $\theta < 1\mathrm{rad}$ 时可作近似 $\sin \theta \approx \theta$ .
\end{problem}
\clearpage


\begin{problem}{2.3.3}
一质量为 $m$ 的质点受到两个力的作用, $\bm{F}_1 = f(r)\bm{e}_r$ , $\bm{F}_2 = -\lambda \bm{v} (\lambda > 0$ , 为常量), $\bm{v}$ 是质点的速度. 若该质点初始时对原点的角动量为 $\bm{J}_0$ , 求以后时刻它对原点的角动量.
\end{problem}
\clearpage


\begin{problem}{2.3.4}
一质量为 m 的粒子, 在半顶角为 $\alpha$ 的锥体内表面无摩擦地滑动.

(1) 找出使该质点绕竖直轴做圆轨道运动在初始条件方面的要求；

(2) 说明这种轨道运动是否稳定,何故?
\end{problem}
\clearpage


\begin{problem}{2.3.5}
质量分别为 $m, 2m, 3m, 4m, 5m$ 的五个质点，用八根长度均为 $l$ 的轻质杆连成图2.71所示的锥体，以角速度 $\omega$ 绕连接 $4m$ 与 $5m$ 的线转动。求锥体对 $O$ 点的角动量的 $x, y, z$ 方向的分量。
\begin{center}
\begin{tikzpicture}[>=Stealth, thick]
  % Coordinate axes
  \coordinate (O) at (0,0);
  \draw[->] (O) -- (0,3.2) node[right] {$z$};
  \draw[->] (O) -- (2.9,0) node[below] {$y$};
  \draw[->] (O) -- (-1.65,-1.14) node[above left=-1pt] {$x$};

  % Pyramid vertices
  \coordinate (M1) at (-2.05,1.22);
  \coordinate (M2) at (-0.75,2.12);
  \coordinate (M3) at (2.35,2.12);
  \coordinate (M4) at (1.05,1.22);

  % Pyramid rods
  \draw (M1) -- (M2) -- (M3) -- (M4) -- cycle;
  \draw (M2) -- (M4);
  \draw (O) -- (M1);
  \draw (O) -- (M2);
  \draw (O) -- (M3);
  \draw (O) -- (M4);

  % Point masses and labels
  \fill (O) circle (2pt) node[left=3pt] {$O$} node[below right=1pt] {$5m$};
  \fill (M1) circle (2pt) node[left=2pt] {$m$};
  \fill (M2) circle (2pt) node[above left=1pt] {$2m$};
  \fill (M3) circle (2pt) node[right=2pt] {$3m$};
  \fill (M4) circle (2pt) node[above=2pt] {$4m$};
\end{tikzpicture}

\end{center}
\end{problem}
\clearpage


\begin{problem}{2.3.6}
一有方向的直线 $L$ ，其方向余弦为 $l, m, n$ ，且通过点 $(a, b, c)$ 。求在 $xy$ 平面上通过原点与 $x$ 轴夹角为 $\alpha$ 、与 $y$ 轴夹角为 $\frac{\pi}{2} - \alpha$ 的单位矢量对有向直线 $L$ 的矩。
\end{problem}
\clearpage


\begin{problem}{2.3.7}
(1) 如果加给 $n$ 个质点的所有力的矢量和等于零，证明：加给这 $n$ 个质点的力矩之矢量和与计算它的参考点（矩心）无关；

(2) 如果 n 个质点的动量之矢量和等于零, 证明: 这 n 个质点的角动量之矢量和与计算它的参考点(矩心)无关.

2.72 所示.
\end{problem}
\clearpage


\begin{problem}{2.3.8}
一个质点在重力作用下，沿一个对竖直轴对称的光滑旋转面上运动。用柱坐标表示的该旋转面的方程为 $\rho = f(z)$ 。如质点在高度 $z_{1}$ 处，速度是水平的、大小为 $v_{1}$ ，在高度 $z_{2}$ 处，速度再次是水平的、大小为 $v_{2}$ 。求 $v_{1}, v_{2}$ ，把它们写成 $z_{1}, z_{2}$ 的函数。
\end{problem}
\clearpage


\begin{problem}{2.3.9}
一个原长为 $a$ 、弹性模量为 $\lambda$ 的弹性绳一端固定在一个光滑水平面上的固定点 $O$ , 另一端系一个质量为 $m$ 的质点, 原先在平面上处于静止, 突然给质点施加一个大小为 $v$ 、方向垂直于绳的速度, 以后运动中, 绳被拉伸到最大长度 $3a$ . 求质点的初速度 $v$ .
\end{problem}
\clearpage


\begin{problem}{2.3.10}
质量为 $m$ 、带电量为 $e$ 的点电荷，在强度为 $g$ 的磁单极子场中运动，磁单极子可以看成无限重、位于原点，点电荷受磁单极子的作用力为 $-ge\frac{\dot{r} \times r}{r^3}$ ，不计重力。

(1) 证明动能 $T=\frac{1}{2}m\dot{r}\cdot\dot{r}$ 是运动积分；

(2) 证明 $L = J + \frac{eg}{r} r$ 是运动积分, 其中 $J = r \times m \dot{r}$ ;

(3) 利用上述 L 是运动积分, 证明这个带电粒子的运动轨道在一个半顶角为 $\alpha$ 的圆锥面上, L 是该圆锥的对称轴, $\alpha$ 满足 $\cos\alpha=\frac{eg}{|L|}$ .

提示:考虑 $r \cdot L$ .
\end{problem}
\clearpage


\begin{problem}{2.3.11}
两个同轴的长半圆筒, 其截面如图 2.74 所示, 分别带正、负电, 使其间电场为 $E = \frac{k}{r} e_r$ , (k 为常量), 一质量为 $m$ 、带电 $-q(q > 0)$ 的粒子, 从左方以垂直半圆筒的轴和垂直半径方向的速度 $v$ 进入两极之间, 如图 2.74 所示. 粒子将在图平面内运动.

(1) 如果粒子在两极间沿圆轨道运动, 其轨道半径 r 多大?

(2) 假定粒子进入场区时, 离轴的距离 $r$ 与速度 $\bm{v}$ 的大小均与 (1) 中相同, 但 $\bm{v}$ 的方向偏离一个小角 $\beta$ . 新轨道与 (1) 中的轨道将再次相交于 $P$ 点, $P$ 点的位置与 $\beta$ 无关, 求出 $P$ 点的位置 (粒子进入场区时仍没有轴向的分速度, 仍在图平面内运动).

(3) 如果把电场改为平行于半圆筒轴的均匀磁场, 粒子沿圆轨道运动, 轨道半径 r 多大?
\begin{center}
\begin{tikzpicture}[>=Stealth, thick]
  \coordinate (C) at (0.6,0);
  \def\Rout{1.95}
  \def\Rin{0.62}
  \def\Rorb{1.15}

  % Outer negatively charged half-cylinder
  \draw (C) ++(-92:\Rout) arc (-92:92:\Rout);
  \foreach \a in {-80,-53,-26,0,26,53,80} {
    \node[font=\footnotesize] at ($(C) + (\a:\Rout+0.20)$) {$-$};
  }

  % Inner positively charged half-cylinder
  \draw (C) ++(-105:\Rin) arc (-105:105:\Rin);
  \foreach \a in {-72,-36,0,36,72} {
    \node[font=\footnotesize] at ($(C) + (\a:\Rin-0.20)$) {$+$};
  }

  % Incident particle and velocity
  \coordinate (S) at (-2.15,\Rorb);
  \coordinate (E) at (0.6,\Rorb);
  \fill (S) circle (1.8pt) node[above=2pt] {$m$} node[below=2pt] {$-q$};
  \draw[->] (S) -- (-1.35,\Rorb) node[midway, above=1pt] {$v$};

  % Circular reference trajectory (dashed)
  \draw[dashed] (-1.35,\Rorb) -- (E)
    arc (90:-90:\Rorb) -- (0.0,-\Rorb);

  % Perturbed trajectory entering at small angle beta
  \coordinate (B0) at (-1.55,0.68);
  \draw[dashed] (B0) -- (E);

  % Angle beta between incident rays
  \coordinate (H) at (-1.35,\Rorb);
  \pic[draw, thin, angle radius=1.55cm, angle eccentricity=1.14, "$\beta$"] {angle = H--E--B0};

  % Perturbed trajectory inside field intersecting at P
  \coordinate (P) at ($(C) + (-37:\Rorb)$);
  \draw[dashed] (E)
    .. controls (1.46,1.32) and (2.06,0.68) .. (P)
    .. controls (1.28,-0.98) and (0.70,-1.28) .. (-0.05,-1.35);
  \node[below right=1pt] at (P) {$P$};
\end{tikzpicture}

\end{center}
\end{problem}
\clearpage


\begin{problem}{2.3.12}
证明: 在两个质量为 $m_{1}, m_{2}$ 的质点按库仑定律相互作用的问题中, 矢量 $L = V \times J + \frac{ar}{r}$ 是运动积分, 其中 $V$ 是相对速度, $r$ 是相对位矢, $J = \mu(r \times V)$ 是与相对运动相联系的角动量, $a$ 是库仑定律中的常系数, $\mu = \frac{m_1 m_2}{m_1 + m_2}$ .

提示:要考虑质点 2 相对于质点 1 的运动.
\end{problem}
\clearpage


\begin{problem}{2.3.13}
一宇宙飞船围绕质量为 $M$ 的恒星作半径为 $r_0$ 的圆周运动, 宇宙飞船的火箭可以迅速点燃, 使飞船速度瞬时改变 $\Delta v$ , 点燃的方位用飞船的速度 $v$ 和 $\Delta v$ 之间的夹角来说明, 如图 2.75 所示, 为了保存燃料, 在连续 $N$ 次发射中, 希望 $\Delta v = \sum_{i=1}^{N} |\Delta v_i|$ 减至最小, 称 $\Delta v$ 为比冲量.

（1）假如我们要利用飞船的火箭逃离这个恒星，且只点燃一次，需要的最小的比冲量多大？在什么方位点燃？

(2) 假如我们要在半径为 $r_{1}(>r_{0})$ 的圆轨道上视察一颗行星, 让火箭再次点燃, 达到该行星轨道所需的最小比冲量多大?

(3) 假如我们要用飞船的火箭使它撞上恒星(假定恒星的半径可以忽略), 对下列两种点火策略计算最小比冲量: (a) 在 $\theta = 180^{\circ}$ 一次迅速点燃; (b) 在 $\theta = 0^{\circ}$ 一次快速点燃, 然后在晚些时候在 $\theta = 180^{\circ}$ 第二次点燃, 为使总的比冲量最小, 选择第二次点火的时间和每次突发的强度.
\begin{center}
\begin{tikzpicture}[>=Stealth, thick]
  % Rocket body tilted slightly up-right
  \begin{scope}[rotate=12]
    \coordinate (C) at (0,0);
    \coordinate (Nose) at (1.15,0);
    % Streamlined fuselage with upper and lower fins
    \draw[fill=white]
      (-1.05,0.18) -- (-1.30,0.42) -- (-0.85,0.28)
      .. controls (-0.25,0.38) and (0.65,0.34) .. (Nose)
      .. controls (0.65,-0.34) and (-0.25,-0.38) .. (-0.85,-0.28)
      -- (-1.30,-0.42) -- (-1.05,-0.18) -- cycle;

    % Exhaust plumes behind nozzle
    \draw (-1.28,0.09) -- (-2.15,0.18);
    \draw (-1.32,0.01) -- (-2.25,-0.02);
    \draw (-1.30,-0.07) -- (-2.12,-0.18);
    \draw (-1.34,-0.16) -- (-2.15,-0.38);

    % Centerline and velocity vector v
    \draw[dashed] (-2.35,0) -- (Nose);
    \draw[->] (Nose) -- (2.35,0) node[above=1pt] {$v$};
    \coordinate (AxisF) at (2.0,0);
  \end{scope}

  % Velocity impulse vector Delta v
  \coordinate (DV) at (1.75,-0.75);
  \draw[->] (C) -- (DV) node[right=1pt] {$\Delta v$};

  % Angle theta between v and Delta v
  \pic[draw, thin, angle radius=1.25cm, angle eccentricity=1.18, "$\theta$"] {angle = DV--C--AxisF};
\end{tikzpicture}

\end{center}
\end{problem}
\clearpage


\section{碰 撞}

\begin{problem}{2.4.1}
一质量为 m 的子弹射入置于光滑水平面上质量为 M 并与劲度系数为 k 的轻弹簧连着的木块后使弹簧最大压缩了 L，求子弹射入前的速度 $v_{0}$ .
\begin{center}
\begin{tikzpicture}[>=Stealth, thick]
  % Horizontal floor and right wall with hatching
  \fill[pattern=north east lines] (0,-0.22) rectangle (5.4,0);
  \fill[pattern=north east lines] (5.18,0) rectangle (5.4,1.65);
  \draw (0,0) -- (5.18,0) -- (5.18,1.65);

  % Bullet of mass m and initial velocity v_0
  \draw (0.35,0.48) -- (0.75,0.58) -- (0.35,0.68) -- cycle;
  \node at (0.52,0.88) {$m$};
  \draw[->] (0.98,0.58) -- (1.68,0.58) node[midway, above=2pt] {$v_0$};

  % Block M on smooth floor
  \draw (1.95,0) rectangle (3.05,1.10);
  \node at (2.50,0.55) {$M$};

  % Spring of stiffness k connected to wall
  \draw (3.05,0.55) -- (3.35,0.55);
  \draw[decorate, decoration={zigzag, segment length=7pt, amplitude=4pt}]
    (3.35,0.55) -- (4.88,0.55);
  \draw (4.88,0.55) -- (5.18,0.55);
  \node at (4.12,0.88) {$k$};
\end{tikzpicture}

\end{center}
\end{problem}
\clearpage


\begin{problem}{2.4.2}
质量为 m、速度为 v 的子弹射向质量为 M 的靶，靶中开有一孔，孔内装了一个劲度系数为 k 的弹簧，如图 2.78 所示。靶在初始时刻处于静止状态，能在无摩擦的水平面上滑动。求子弹射入靶后弹簧的最大压缩距离 $\Delta x$ 。
\begin{center}
\begin{tikzpicture}[>=Stealth, thick]
  % Smooth horizontal floor with hatching
  \fill[pattern=north east lines] (0,-0.22) rectangle (5.4,0);
  \draw (0,0) -- (5.4,0);

  % Target block M with horizontal cavity open on the right
  \draw (0.75,0) -- (1.95,0) -- (1.95,0.24) -- (1.05,0.24)
    -- (1.05,0.78) -- (1.95,0.78) -- (1.95,1.02) -- (0.75,1.02) -- cycle;
  \node at (1.35,1.32) {$M$};

  % Coiled spring inside the cavity
  \draw (1.05,0.51) -- (1.20,0.51);
  \draw[decorate, decoration={coil, aspect=0.45, segment length=4.8pt, amplitude=3.5pt}]
    (1.20,0.51) -- (1.82,0.51);

  % Incoming bullet m moving left with velocity v
  \draw[->] (3.95,0.51) -- (3.20,0.51) node[midway, above=2pt] {$v$};
  \draw (4.82,0.41) -- (4.82,0.61)
    .. controls (4.52,0.62) and (4.35,0.57) .. (4.22,0.51)
    .. controls (4.35,0.45) and (4.52,0.40) .. (4.82,0.41) -- cycle;
  \node at (4.56,0.85) {$m$};
\end{tikzpicture}

\end{center}
\end{problem}
\clearpage


\begin{problem}{2.4.3}
一辆质量为 m 的车以速度 v 驶向另一辆原来静止的质量为 3m 的车，在两车正碰时弹簧被压缩，如图 2.79 所示.

(1) 若机械能是守恒的,那么在弹簧压缩最大的瞬间,质量为 3m 的那辆车的速度多大?

(2) 若机械能不守恒的话, 结果是否不一样?

(3) 若机械能守恒, 经过较长的一段时间后, 较重的车的末速度多大?

(4) 对于完全非弹性碰撞的情形,较重的车的末速度多大?
\begin{center}
\begin{tikzpicture}[>=Stealth, thick]
  % Smooth horizontal floor with hatching
  \fill[pattern=north east lines] (0,-0.22) rectangle (5.6,0);
  \draw (0,0) -- (5.6,0);

  % Left cart of mass m
  \draw (0.95,0.18) circle (0.18);
  \draw (1.65,0.18) circle (0.18);
  \draw (0.72,0.36) rectangle (1.88,1.10);
  \node at (1.30,0.73) {$m$};
  \draw[->] (0.95,1.32) -- (1.65,1.32) node[midway, above=1pt] {$v$};

  % Right cart of mass 3m with attached spring
  \draw (4.00,0.20) circle (0.20);
  \draw (4.85,0.20) circle (0.20);
  \draw (3.60,0.40) rectangle (5.25,1.25);
  \node at (4.42,0.82) {$3m$};

  % Coiled spring on left side of 3m cart
  \draw (3.60,0.82) -- (3.45,0.82);
  \draw[decorate, decoration={coil, aspect=0.45, segment length=4.8pt, amplitude=3.8pt}]
    (3.45,0.82) -- (2.82,0.82);
\end{tikzpicture}

\end{center}
\end{problem}
\clearpage


\begin{problem}{2.4.6}
一质量为 $M$ 的车上装有一根杆，用一细绳将质量为 $\mu$ 的球挂在杆上 $P$ 点，车和球的初速度为 $v$ ，这辆车撞到另一辆质量为 $m$ 的原来静止的车上并且和它粘在一起.若细绳的长度为 $R$ ，试证明能使球绕 $P$ 点转圆圈的最小初速度为 $v = \left(\frac{m + M}{m}\right)\sqrt{5gR}$ ，忽

略摩擦并假定 $M, m \gg \mu.$
\begin{center}
\begin{tikzpicture}[>=Stealth, thick]
  % Smooth horizontal floor with hatching
  \fill[pattern=north east lines] (0,-0.22) rectangle (5.6,0);
  \draw (0,0) -- (5.6,0);

  % Left cart of mass M with pole and pendulum
  \draw (0.95,0.20) circle (0.20);
  \draw (1.80,0.20) circle (0.20);
  \draw (0.55,0.40) rectangle (2.20,1.20);
  \node at (1.375,0.80) {$M$};

  % Vertical pole and pendulum bob mu
  \coordinate (P) at (1.32,2.35);
  \coordinate (Bob) at (1.68,1.56);
  \draw (1.32,1.20) -- (P);
  \node[above=1pt] at (P) {$P$};
  \draw (P) -- (Bob);
  \fill (Bob) circle (2.2pt) node[right=2pt] {$\mu$};

  % Velocity arrow v
  \draw[->] (2.35,0.80) -- (2.98,0.80) node[midway, above=1pt] {$v$};

  % Right stationary cart of mass m
  \draw (3.85,0.20) circle (0.20);
  \draw (4.70,0.20) circle (0.20);
  \draw (3.45,0.40) rectangle (5.10,1.20);
  \node at (4.275,0.80) {$m$};
\end{tikzpicture}

\end{center}
\end{problem}
\clearpage


\begin{problem}{2.4.7}
一弹子球从阶梯上弹下，在每一台阶都落在同一位置，都弹起相同的高度，每个台阶的高度和宽度相等为 $l$ ，每次碰撞恢复系数为 $\bm{e}$ .求弹子球所必需的水平速度 $\bm{v}_{h}$ 和每次弹起的高度 $H$
\end{problem}
\clearpage


\begin{problem}{2.4.8}
一质量为 m、荷电为 q、初速为 v 的粒子，和另一个处于静止的全同粒子发生正碰，求：

(1) 两个粒子最接近时的距离；

(2) 最接近的瞬时两个粒子的速度；

(3) 两个粒子的末速度.
\end{problem}
\clearpage


\begin{problem}{2.4.9}
小球 1 从离碗底高 h 处静止滑下,与静置于碗底的完全相同的小球 2 发生非

弹性碰撞,若恢复系数为 e,求第一次碰撞后两球达到的高度(忽略一切摩擦力).
\end{problem}
\clearpage


\begin{problem}{2.4.10}
质量为 $m_{1}$ 和 $m_{2}$ 的两个小球悬挂在长度分别为 $l_{1}$ 和 $l_{2}$ 的不可伸长的轻绳下，两球恰好相切。在两线所在的平面内把第一个小球拉到与铅垂线成 $\alpha$ 角的位置，静止后再放开，摆下与静止的第二个球发生弹性碰撞。求第一次碰撞后两球偏离铅垂线的最大角度 $\alpha_{1}$ 和 $\alpha_{2}$ 。
\end{problem}
\clearpage


\begin{problem}{2.4.11}
一个运动粒子与另一个静止的全同粒子发生弹性碰撞,试证明:碰后两粒子的运动方向相互垂直.
\end{problem}
\clearpage


\begin{problem}{2.4.12}
两个完全相同的光滑小球,一个静止,另一个以 $v_{0}$ 的速度与之发生非弹性碰撞.测得原静止的小球以 $\frac{3}{8}v_{0}$ 的速率向着偏离 $v_{0}$ 方向60°的方向运动.求:

(1) 恢复系数 e,

(2) 原来运动的小球的散射角.
\begin{center}
\begin{tikzpicture}[>=Stealth, thick]
  % Coordinate axes
  \coordinate (O) at (0,0);
  \draw[->] (O) -- (0,2.55) node[left=1pt] {$y$};
  \draw[->] (O) -- (4.35,0) node[right=1pt] {$x$};

  % Two identical touching spheres
  \coordinate (C1) at (1.50,1.15);
  \coordinate (C2) at (3.20,1.15);
  \draw (C1) circle (0.85);
  \draw (C2) circle (0.85);
  \fill (C1) circle (1.4pt);
  \fill (C2) circle (1.4pt);
  \node at (0.68,2.08) {$1$};
  \node at (3.20,2.22) {$2$};

  % Velocity components and initial velocity on sphere 1
  \coordinate (V1x) at (2.22,1.15);
  \coordinate (V1y) at (1.50,2.60);
  \coordinate (V0)  at (2.34,2.60);
  \draw[->] (C1) -- (V1x) node[midway, below=2pt] {$v_{1x}$};
  \draw[->] (C1) -- (V1y) node[left=1pt] {$v_{1y}$};
  \draw[->] (C1) -- (V0)  node[above right=-1pt] {$v_0$};
  \draw[dashed] (V1y) -- (V0);

  % 60 degree angle
  \pic[draw, thin, angle radius=0.42cm, angle eccentricity=1.48, "$60^\circ$"] {angle = V1x--C1--V0};

  % Velocity of sphere 2
  \draw[->] (C2) -- (4.55,1.15) node[right=1pt] {$v_2=\frac{3}{8}v_0$};
\end{tikzpicture}

\end{center}
\end{problem}
\clearpage


\begin{problem}{2.4.13}
一个质量为 $m_{1}$ 以 $\bm{v}_{10}$ 的速度运动的粒子1，与一个质量为 $m_{2}$ 的静止粒子2发生弹性斜碰撞，碰撞后，在质心平动参考系中，粒子1的散射角为 $\alpha_{\mathrm{c}}$ 。试证：在实验室参考系中两粒子的速度与 $\bm{v}_{10}$ 的夹角 $\alpha, \beta$ 满足下列关系：

\[
\tan \alpha = \frac {m _ {2} \sin \alpha_ {\mathrm{c}}}{m _ {1} + m _ {2} \cos \alpha_ {\mathrm{c}}}, \quad \tan \beta = \frac {\sin \alpha_ {\mathrm{c}}}{1 - \cos \alpha_ {\mathrm{c}}}
\]
\end{problem}
\clearpage


\begin{problem}{2.4.14}
一质量为 $m$ 的质子以 $v_{0}(v_{0} \ll c)$ 的速度去撞击静止的质量为 $4m$ 的氦核，用实验室参考系，质子以 $\frac{1}{2} v_{0}$ 的速率和 $30^{\circ}$ 的角度散射。求：

（1）用实验室参考系，撞击后氦核的速率及运动方向；

(2) 用质心平动参考系, 质子的散射速度.
\begin{center}
\begin{tikzpicture}[>=Stealth, thick]
  % Coordinate axes (x right, y down)
  \coordinate (O) at (0,1.45);
  \draw[->] (O) -- (1.05,1.45) node[right=1pt] {$x$};
  \draw[->] (O) -- (0,0.35) node[below=1pt] {$y$};

  % Proton (top) and its scattered velocity
  \coordinate (Cp) at (2.45,1.45);
  \draw (Cp) circle (0.30);
  \node at (1.92,1.88) {质子};
  \coordinate (Href) at (4.15,1.45);
  \coordinate (Vp) at ($(Cp) + (30:1.90)$);
  \draw[dashed] (Cp) -- (Href);
  \draw[->] (Cp) -- (Vp) node[above right=-1pt] {$\frac{1}{2}v_0$};
  \pic[draw, thin, angle radius=0.82cm, angle eccentricity=1.28, "$30^\circ$"] {angle = Href--Cp--Vp};

  % Helium nucleus (bottom) and its recoil velocity components
  \coordinate (Che) at (3.00,0.72);
  \draw (Che) circle (0.45);
  \node at (2.05,0.68) {氦核};
  \coordinate (V2x) at (4.50,0.72);
  \coordinate (V2y) at (3.00,-0.12);
  \coordinate (V2)  at (4.50,-0.12);
  \draw[->] (Che) -- (V2x) node[right=1pt] {$v_{2x}$};
  \draw[->] (Che) -- (V2y) node[below left=-1pt] {$v_{2y}$};
  \draw[->] (Che) -- (V2);
  \draw[dashed] (V2x) -- (V2) -- (V2y);
  \pic[draw, thin, angle radius=0.72cm, angle eccentricity=1.28, "$\alpha$"] {angle = V2--Che--V2x};
\end{tikzpicture}

\end{center}
\end{problem}
\clearpage

